% Báo cáo cuối kỳ — Đề tài 9.
% Biên dịch: tectonic docs/report/report.tex (hoặc XeLaTeX ba lượt để ổn định mục lục).
% Nguồn nội dung: docs/report/sections/*.tex
\documentclass[11pt,a4paper]{article}
\usepackage[margin=22mm,headheight=14pt]{geometry}
\usepackage{fontspec}
\setmainfont{texgyretermes}[Extension=.otf,UprightFont=*-regular,BoldFont=*-bold,ItalicFont=*-italic,BoldItalicFont=*-bolditalic]
% Không đặt \setmonofont để tránh phụ thuộc thêm font; \texttt dùng font mặc định.
\usepackage{amsmath,amssymb}
\usepackage{booktabs}
\usepackage{longtable}
\usepackage{array}
\usepackage{fancyhdr}
\usepackage{enumitem}
\usepackage[unicode,hidelinks]{hyperref}
\usepackage{bookmark}

\hypersetup{pdftitle={Đề tài 9 — Kiến trúc Transformer và tinh chỉnh mô hình ngôn ngữ},pdflang={vi-VN}}
\setlength{\parindent}{0pt}
\setlength{\parskip}{5pt plus 1pt minus 1pt}
\setlength{\emergencystretch}{2em}
\linespread{1.08}
\raggedbottom
\renewcommand{\contentsname}{Mục lục}
\renewcommand{\refname}{Tài liệu tham khảo}
\renewcommand{\headrulewidth}{0.3pt}
\pagestyle{fancy}
\fancyhf{}
\fancyhead[L]{\small Đề tài 9 — Báo cáo cuối kỳ}
\fancyhead[R]{\small Bản nháp chưa nhóm duyệt}
\fancyfoot[C]{\thepage}

% Ký hiệu dùng chung với docs/theory/part-b.tex
\DeclareMathOperator{\Var}{Var}
\DeclareMathOperator{\Cov}{Cov}
\DeclareMathOperator{\rank}{rank}
\DeclareMathOperator{\im}{im}
\DeclareMathOperator{\diag}{diag}
\DeclareMathOperator{\softmax}{softmax}
\DeclareMathOperator{\KL}{KL}
\DeclareMathOperator{\supp}{supp}
\DeclareMathOperator*{\argmin}{arg\,min}
\DeclareMathOperator*{\argmax}{arg\,max}
\newcommand{\E}{\mathbb{E}}
\newcommand{\R}{\mathbb{R}}
\newcommand{\PE}{\mathrm{PE}}
\newcommand{\Attn}{\operatorname{Attn}}
\newcommand{\PPL}{\mathrm{PPL}}
\newcommand{\pref}{\pi_{\mathrm{ref}}}
\newcommand{\pdata}{p_{\mathrm{data}}}
\newcommand{\dmodel}{d_{\mathrm{model}}}

% Nhãn trạng thái: dùng để phân biệt việc đã làm với việc chưa làm.
\newcommand{\dalam}{\textbf{[đã thực hiện]}}
\newcommand{\chualam}{\textbf{[chưa thực hiện]}}
\newcommand{\motphan}{\textbf{[một phần]}}

\begin{document}

%====================================================================
\begin{titlepage}
\centering
\vspace*{18mm}

{\large BÁO CÁO BÀI TẬP LỚN\par}
{\large HỌC MÁY CƠ BẢN\par}

\vspace{20mm}
{\LARGE\bfseries Đề tài 9\par}
\vspace{4mm}
{\LARGE\bfseries Kiến trúc Transformer và tinh chỉnh\\[2mm] mô hình ngôn ngữ\par}

\vspace{8mm}
{\large\itshape Ứng dụng: trợ lý hỏi đáp thủ tục hộ tịch\\ và chứng thực tại Hà Nội\par}

\vspace{24mm}
\begin{tabular}{rl}
\textbf{Nhóm thực hiện:} & Cao Văn Hà (trưởng nhóm)\\
 & Hà Mạnh Thế\\
 & Trần Duy Hưng\\[4mm]
\textbf{Đề gốc:} & Bản ngày 12/11/2026, Đề tài 9, tr.~18--19\\[4mm]
\textbf{Hạn hoàn thành:} & 23h59 ngày 01/12/2026\\
\textbf{Vấn đáp:} & 03/12/2026\\
\end{tabular}

\vfill
\begin{minipage}{0.86\textwidth}
\small
\textbf{Trạng thái tài liệu.} Đây là bản nháp có hỗ trợ AI, chưa được nhóm xác nhận là
bản nộp. Mỗi mục đều ghi nhãn \dalam, \motphan{} hoặc \chualam{} theo bằng chứng
thực tế trong kho mã. Những phần chưa có kết quả thực nghiệm không được trình bày
như đã hoàn thành. Phạm vi hỗ trợ AI khai báo tại Phụ lục~\ref{sec:phu-luc-ai}.
\end{minipage}

\vspace{10mm}
\end{titlepage}

%====================================================================
\pagenumbering{roman}
\tableofcontents
\clearpage
\section*{Tóm tắt}
\addcontentsline{toc}{section}{Tóm tắt}

Báo cáo trình bày bốn phần bắt buộc của Đề tài 9 \cite{de}: cài đặt từ đầu các
thành phần Transformer và hai hàm mục tiêu tinh chỉnh (Phần~A), chứng minh sáu
mệnh đề lý thuyết nền (Phần~B), bài tập tính toán nối lý thuyết với cài đặt
(Phần~C), và ứng dụng tinh chỉnh một mô hình ngôn ngữ thành trợ lý hỏi đáp thủ
tục hộ tịch -- chứng thực tại Hà Nội (Phần~D).

\textbf{Phần A \dalam.} Nhóm cài đặt bằng PyTorch, không dùng lớp attention hay
vòng huấn luyện dựng sẵn: chú ý tích vô hướng có chia tỷ lệ, chú ý nhiều đầu,
mặt nạ nhân quả, mã hóa vị trí hình sin, một khối decoder pre-LN, lớp LoRA có
hợp nhất trọng số, và hàm mất mát DPO. Mỗi cài đặt được đối chiếu với thư viện
chuẩn trên cùng đầu vào. Mô hình ngôn ngữ ký tự huấn luyện trên bản
\emph{Truyện Kiều} Trương Vĩnh Ký 1911 đạt perplexity validation
$\mathbf{10{,}1948 \pm 0{,}0628}$ (trung bình $\pm$ độ lệch chuẩn mẫu của 5 seed
độc lập).

\textbf{Phần B \dalam.} Sáu chứng minh đầy đủ giả thiết: phương sai tích vô hướng
bằng $d_k$ và vai trò hệ số $1/\sqrt{d_k}$; tính đẳng biến hoán vị của
self-attention cùng phản ví dụ cho thấy mặt nạ nhân quả phá tính chất này; biểu
diễn độ dịch vị trí bằng ma trận quay khối, kèm chứng minh mệnh đề \emph{thất bại}
ở số chiều lẻ; bốn tiểu mệnh đề LoRA; nghiệm duy nhất của bài toán reward có phạt
KL và dẫn xuất mất mát DPO qua Bradley--Terry; tương đương giữa cực tiểu
cross-entropy và cực tiểu $\KL(\pdata\|p_\theta)$ cùng quan hệ với perplexity.

\textbf{Phần C \dalam.} Bốn bài tập được giải và đối chiếu số với mã Phần~A: một
đầu attention nhân quả trên chuỗi 3 token; đếm $\mathbf{33{.}472}$ tham số và
$\mathbf{6{.}291{.}456}$ phép nhân--cộng cho một khối decoder theo cấu hình A.2 ở
độ dài 128, kèm phân tích hạng tử $O(n^2)$; số tham số LoRA ở Q/V của
Qwen2.5-1.5B-Instruct là $\mathbf{544{.}768}$, $\mathbf{1{.}089{.}536}$,
$\mathbf{2{.}179{.}072}$ cho $r=4,8,16$ cùng ước lượng bộ nhớ khi kết hợp QLoRA
4-bit; tính mất mát và gradient DPO, phân tích ảnh hưởng của $\beta$.

\textbf{Phần D \motphan.} Nhóm đã rà năm họ thủ tục từ nguồn chính văn, dựng quy
trình sinh QA bán tự động có cổng phát hành, và \textbf{chạy xong một lượt SFT
pilot thật} trên Kaggle: lô 24 QA đã duyệt 24/24, QLoRA $r=8$ ở Q/V, validation
PPL giảm $4{,}53\to3{,}46$ sau 2 epoch, peak allocated $2{,}402$ GiB, resume tái
lập đúng chỉ số. Tập chính đã lên \textbf{400 QA ứng viên} nhưng chưa phát hành.
Hỏi thử adapter pilot cho thấy hai lỗi thật: sai thời hạn, và \emph{không từ chối}
câu ngoài phạm vi. Chưa thu cặp so sánh, chưa huấn luyện DPO, chưa đánh giá ba mốc.
Mục~\ref{sec:phan-d} ghi rõ từng ngưỡng và khoảng cách còn lại.

\textbf{Kết luận trung thực.} Ba phần lý thuyết và cài đặt (A, B, C) có bằng chứng
kiểm chứng được trong kho mã. Phần ứng dụng còn thiếu các hạng mục phụ thuộc vào
con người thật: người khảo sát có văn bản đồng ý cho tập test, hai người chấm độc
lập cho tối thiểu 500 cặp ưu tiên, và người đánh giá mù cho so sánh cuối. Báo cáo
không trình bày số liệu giả định như kết quả đo được.

\clearpage
\pagenumbering{arabic}

\section{Giới thiệu}
\label{sec:gioi-thieu}

\subsection{Bài toán và động cơ}

Thủ tục hộ tịch và chứng thực là nhóm thủ tục hành chính mà người dân tiếp xúc
nhiều nhất: khai sinh, khai tử, kết hôn, cấp trích lục, chứng thực bản sao và
chứng thực chữ ký. Thông tin về trình tự, thành phần hồ sơ, thời hạn giải quyết
và lệ phí đều được công bố công khai, nhưng phân tán trên nhiều văn bản ở nhiều
cấp hiệu lực khác nhau, và thay đổi theo thời gian. Người dân thường không biết
điều kiện áp dụng nào tương ứng với trường hợp của mình.

Một trợ lý hỏi đáp có thể giảm chi phí tra cứu, nhưng chỉ khi nó thỏa ba điều
kiện. Thứ nhất, câu trả lời phải nêu đúng điều kiện áp dụng, không chỉ nêu kết
luận chung. Thứ hai, câu trả lời phải \emph{dẫn nguồn} để người dùng tự kiểm tra,
vì giá trị pháp lý thuộc về văn bản gốc chứ không thuộc về mô hình. Thứ ba, khi
câu hỏi ngoài phạm vi hoặc chưa đủ căn cứ, trợ lý phải nói rõ thay vì đoán. Điều
kiện thứ ba quan trọng hơn điều thứ nhất về mặt rủi ro: một câu trả lời sai nhưng
tự tin gây hại nhiều hơn một lời từ chối.

Ba điều kiện này định hình toàn bộ thiết kế Phần~D. Chúng không đạt được bằng cách
huấn luyện mô hình cho đúng nhiều hơn, mà cần một tín hiệu học phân biệt được
``đúng và có dẫn nguồn'' với ``trôi chảy nhưng không kiểm được''. Đó chính là lý do
đề tài ghép SFT với tối ưu hóa theo ưu tiên: SFT dạy định dạng và nội dung, còn DPO
dạy thứ tự ưu tiên giữa hai câu trả lời mà con người đã xếp hạng.

\subsection{Phạm vi bắt buộc theo đề}

Đề tài 9 \cite[tr.~18--19]{de} yêu cầu bốn phần trên cùng một chủ đề, với trọng số
điểm như sau.

\begin{center}
\begin{tabular}{llr}
\toprule
\textbf{Phần} & \textbf{Nội dung} & \textbf{Điểm}\\
\midrule
A & Cài đặt từ đầu, có dẫn xuất và đối chiếu thư viện chuẩn & 20\\
B & Chứng minh chặt chẽ, đầy đủ giả thiết & 25\\
C & Bài tập tính toán, phân tích, thực nghiệm có kiểm soát & 15\\
D & Ứng dụng: thu thập/gán nhãn, tinh chỉnh, đánh giá & 30\\
\midrule
 & Báo cáo, tái lập, bảo vệ & 10\\
\bottomrule
\end{tabular}
\end{center}

Phần~D có trọng số lớn nhất. Đề cũng ràng buộc rằng ít nhất một cài đặt của
Phần~A phải được dùng trực tiếp hoặc so sánh trong quy trình Phần~D; Phần~D được
phép dùng thư viện chuẩn nhưng điều đó không thay thế nghĩa vụ tự cài đặt ở
Phần~A.

Hai phần lý thuyết B và C phải dựa trên năm bài báo trong danh mục tài liệu tham
khảo của đề: \cite{attention}, \cite{lora}, \cite{qlora}, \cite{rlhf},
\cite{dpo}. Báo cáo ghi rõ mục và phương trình được dùng ở từng chỗ.

\subsection{Đầu vào và đầu ra của sản phẩm}

Sản phẩm Phần~D là một trợ lý hỏi đáp.

\begin{description}[leftmargin=2em,style=nextline]
\item[Đầu vào] Một câu hỏi tiếng Việt về thủ tục hộ tịch hoặc chứng thực tại
Hà Nội, do người dùng tự diễn đạt. Không giả định người dùng biết tên thủ tục,
số hiệu văn bản hay thuật ngữ hành chính.
\item[Đầu ra mong muốn] Câu trả lời nêu đúng thông tin và điều kiện áp dụng, kèm
dẫn nguồn để người dùng kiểm tra. Nếu câu hỏi ngoài phạm vi hộ tịch -- chứng thực
tại Hà Nội, hoặc chưa đủ căn cứ trong dữ liệu đã thu thập, trợ lý nêu rõ điều đó
thay vì đoán.
\end{description}

Hai điều kiện ``dẫn nguồn'' và ``từ chối khi chưa đủ căn cứ'' là tiêu chí đánh giá
chính, không phải tính năng phụ. Chúng xuất hiện trong rubric chấm cặp ưu tiên
(Mục~\ref{sec:phan-d}) và trong giao thức đánh giá cuối.

\subsection{Cấu trúc báo cáo và quy ước trình bày}

Mục~\ref{sec:phan-a} trình bày Phần~A theo trình tự dẫn xuất $\to$ cài đặt $\to$
đối chiếu. Mục~\ref{sec:phan-b} là sáu chứng minh. Mục~\ref{sec:phan-c} là bốn
bài tập tính toán. Mục~\ref{sec:phan-d} là Phần~D, trình bày theo trạng thái thực
tế. Mục~\ref{sec:tai-lap} ghi điều kiện tái lập. Mục~\ref{sec:dao-duc} ghi các
ràng buộc đạo đức và pháp lý. Mục~\ref{sec:gioi-han} nêu giới hạn và hướng mở
rộng. Phụ lục~\ref{sec:phu-luc-ai} khai báo phạm vi dùng AI tạo sinh.

Ba quy ước áp dụng xuyên suốt:

\begin{enumerate}[leftmargin=2em]
\item \textbf{Mọi logarit là logarit tự nhiên}, kể cả trong perplexity và trong
các log-tỷ số của DPO.
\item \textbf{Nhãn trạng thái.} Mỗi hạng mục được gắn \dalam, \motphan{} hoặc
\chualam{} theo bằng chứng trong kho mã. Nhãn không phải dự báo; một hạng mục đã
lập kế hoạch nhưng chưa chạy vẫn là \chualam.
\item \textbf{Phân biệt số đo với số minh họa.} Những ví dụ số dùng để minh họa
phép tính được ghi rõ là minh họa. Chúng không được trình bày như dữ liệu đề chính
thức hay như kết quả thực nghiệm.
\end{enumerate}

\subsection{Phân công trong nhóm}

Nhóm không chia cứng bốn phần cho ba thành viên. Cao Văn Hà thực hiện bản đầu
xuyên suốt A--D, tổng hợp và nộp bài. Hà Mạnh Thế và Trần Duy Hưng luân phiên rà
soát từng đầu việc: chạy lại một phép kiểm mã, tự giải lại một bước chứng minh
hoặc tính toán, hoặc đối chiếu một bản ghi QA với URL và đúng mục trong nguồn.
Cách này nhằm bảo đảm cả ba thành viên đều giải thích được bất kỳ đoạn mã hoặc
chứng minh nào khi vấn đáp, theo yêu cầu phần~II của đề.

Riêng bước tạo nhãn ưu tiên cho Phần~D có quy trình khác và không luân phiên: Hà
Mạnh Thế và Trần Duy Hưng \emph{đều} chấm từng cặp một cách độc lập, không chia
đôi số cặp; Cao Văn Hà chỉ phân xử \emph{sau} khi nhãn độc lập đã bị khóa và mức
đồng thuận đã được đo. Lý do là yêu cầu của đề về hai người đánh giá độc lập chỉ
có nghĩa khi hai nhãn thực sự được tạo độc lập trên cùng một cặp.

\clearpage
\section{Phần A --- Cài đặt từ đầu}
\label{sec:phan-a}

Đề yêu cầu Phần~A phải \emph{dẫn xuất từ bài toán tối ưu hoặc mô hình xác suất},
cài đặt từ đầu bằng NumPy hoặc PyTorch mà không dùng hàm huấn luyện dựng sẵn, rồi
đối chiếu với thư viện chuẩn \cite[tr.~1, mục I]{de}. Mục này theo đúng ba bước đó
cho từng hạng mục. Thư viện chuẩn chỉ xuất hiện ở phía đối chiếu, không xuất hiện
trong đường tính toán của nhóm.

Mã dùng lại nằm trong \texttt{src/}; notebook điều khiển và giải thích nằm trong
\texttt{notebooks/}. Toàn bộ phép kiểm ở mục này chạy trên CPU cục bộ.

%--------------------------------------------------------------------
\subsection{A.1 --- Các thành phần của khối decoder}
\label{sec:a1}

\subsubsection{Dẫn xuất: attention như một phép lấy trung bình có trọng số}

Xuất phát từ bài toán truy xuất mềm. Cho một truy vấn $q\in\R^{d_k}$ và $n$ cặp
khóa--giá trị $(k_j,v_j)$, ta muốn lấy ra giá trị ứng với khóa ``giống'' truy vấn
nhất. Phép lấy cứng $v_{\argmax_j\langle q,k_j\rangle}$ không khả vi, nên thay bằng
một phân phối trên các khóa rồi lấy kỳ vọng. Chọn họ phân phối Gibbs với thế năng
là tích vô hướng:
\[
  \alpha_j=\frac{\exp(\langle q,k_j\rangle/T)}{\sum_{a=1}^{n}\exp(\langle q,k_a\rangle/T)},
  \qquad \text{đầu ra} = \sum_{j=1}^{n}\alpha_j v_j .
\]
Đây chính là nghiệm của bài toán tối ưu có phạt entropy
\[
  \max_{\alpha\in\Delta^{n-1}}\ \sum_j \alpha_j\langle q,k_j\rangle
    + T\Bigl(-\sum_j\alpha_j\log\alpha_j\Bigr),
\]
với $\Delta^{n-1}$ là simplex xác suất; nghiệm là softmax của $\langle q,k_j\rangle/T$.
Cấu trúc này lặp lại ở Mục~\ref{sec:b5} với $T=\beta$ và thế năng là reward, nên
attention và DPO dùng cùng một dạng nghiệm Gibbs.

Chọn $T=\sqrt{d_k}$ không phải tùy ý. Dưới giả thiết các thành phần của $q,k$ độc
lập, kỳ vọng $0$, phương sai $1$, ta chứng minh ở Mục~\ref{sec:b1} rằng
$\Var(q^\top k)=d_k$; chia cho $\sqrt{d_k}$ đưa phương sai mỗi điểm về $1$ và thu
nhỏ mọi chênh lệch điểm theo cùng hệ số, nhờ đó giảm bão hòa softmax
\cite[§3.2.1]{attention}.

Viết theo ma trận cho cả $n$ truy vấn:
\begin{equation}
  \Attn(Q,K,V)=\softmax\!\left(\frac{QK^\top}{\sqrt{d_k}}\right)V,
  \label{eq:sdpa}
\end{equation}
đúng phương trình (1) của \cite{attention}.

\subsubsection{Mặt nạ nhân quả như một ràng buộc hỗ trợ}

Mô hình ngôn ngữ nhân quả phân rã xác suất chuỗi theo quy tắc dây chuyền:
\begin{equation}
  p_\theta(y_1,\ldots,y_T)=\prod_{t=1}^{T}p_\theta(y_t\mid y_{<t}).
  \label{eq:chain}
\end{equation}
Để \eqref{eq:chain} đúng, biểu diễn tại vị trí $t$ không được phụ thuộc
$y_{>t}$. Ràng buộc này áp vào \eqref{eq:sdpa} bằng cách buộc $\alpha_{tj}=0$ với
$j>t$, tương đương cộng $-\infty$ vào điểm số trước softmax. Vi phạm ràng buộc này
làm \eqref{eq:chain} sai, nên hàm mất mát không còn là log-hợp lý của mô hình nào
cả --- đây là lý do mặt nạ là vấn đề đúng/sai, không phải lựa chọn.

Cài đặt đặt mặt nạ dưới dạng boolean ``bị chặn'', giữ nguyên đường chéo:

\begin{quote}\small\ttfamily
def causal\_mask(length, *, device=None):\\
\hspace*{2em}return torch.ones(length, length, dtype=torch.bool, device=device).triu\_(1)
\end{quote}

Hàm \texttt{scaled\_dot\_product\_attention} nhận mặt nạ này và dùng
\texttt{masked\_fill\_} với $-\infty$ trước \texttt{softmax}. Mỗi truy vấn luôn còn
ít nhất một khóa hợp lệ (chính nó), nên mẫu số softmax dương và không sinh
\texttt{NaN}. Tham số \texttt{blocked=None} nghĩa là \emph{không che}, không phải
chế độ nhân quả --- phân biệt này được ghi trong docstring để tránh dùng sai.

\subsubsection{Chú ý nhiều đầu}

Một đầu attention chỉ tạo được một phân phối trọng số trên mỗi vị trí. Để mô hình
dùng nhiều quan hệ song song, chia $\dmodel$ thành $h$ đầu, mỗi đầu có
$d_h=\dmodel/h$, tính attention riêng rồi ghép và chiếu
\cite[§3.2.2]{attention}:
\[
  \mathrm{MHA}(X)=\mathrm{Concat}(H_1,\ldots,H_h)W_O,\qquad
  H_i=\Attn(XW_Q^{(i)},XW_K^{(i)},XW_V^{(i)}).
\]
Vì $h\cdot d_h=\dmodel$, tổng chi phí các phép chiếu không đổi so với một đầu có
cùng $\dmodel$ --- một điểm thường bị hỏi ở vấn đáp và được dùng lại ở
Mục~\ref{sec:c2}.

Cài đặt chiếu Q/K/V \emph{một lần} bằng \texttt{nn.Linear(d\_model, 3*d\_model)} rồi
tách, thay vì ba lớp riêng. Phép biến đổi hình dạng giữ đúng thứ tự trục:

\begin{quote}\small\ttfamily
q, k, v = self.qkv(x).view(batch, length, 3, self.n\_heads, self.head\_dim)\\
\hspace*{6em}.permute(2, 0, 3, 1, 4).unbind(0)
\end{quote}

Cho \texttt{[B,T,3,H,Dh]} $\to$ ba tensor \texttt{[B,H,T,Dh]}. Trục đầu và trục
token không bị trộn; nếu trộn, mỗi đầu sẽ nhận một lát đặc trưng sai và kết quả vẫn
có hình dạng đúng nhưng sai về ngữ nghĩa. Đây là lỗi im lặng nên notebook A.1 kiểm
riêng bằng cách so với \texttt{nn.MultiheadAttention} dùng \emph{cùng trọng số}.

\subsubsection{Mã hóa vị trí hình sin}

Self-attention không có mã hóa vị trí là đẳng biến hoán vị (chứng minh ở
Mục~\ref{sec:b2}), nên không phân biệt được thứ tự token. Transformer gốc cộng một
bảng cố định \cite[§3.5]{attention}:
\[
  \PE(p)_{2i}=\sin\!\left(p\,\omega_i\right),\quad
  \PE(p)_{2i+1}=\cos\!\left(p\,\omega_i\right),\quad
  \omega_i=10000^{-2i/\dmodel}.
\]
Tần số được tính trong miền log để tránh tràn số:

\begin{quote}\small\ttfamily
freq = torch.exp(torch.arange(0, d\_model, 2, dtype=torch.float32)\\
\hspace*{6em}* (-math.log(10000.0) / d\_model))
\end{quote}

Bảng được đăng ký bằng \texttt{register\_buffer(\ldots, persistent=False)}: nó là
hằng số, không phải tham số học, và không cần lưu vào \texttt{state\_dict} vì tái
tạo được từ cấu hình. PE được cộng \emph{đúng một lần} trước chồng block.

Cài đặt chấp nhận $\dmodel$ lẻ; khi đó tọa độ $\sin$ cuối không có $\cos$ đi kèm.
Mục~\ref{sec:b3} chứng minh rằng \emph{tính chất dịch vị trí không còn đúng} trong
trường hợp đó. Thí nghiệm A.2 dùng $\dmodel=64$ nên không chạm phải giới hạn này,
nhưng giới hạn được ghi rõ thay vì bỏ qua.

\subsubsection{Khối decoder pre-LN}

Khối ghép hai nhánh residual, chuẩn hóa \emph{trước} mỗi nhánh:
\[
  x \leftarrow x + \mathrm{Attn}(\mathrm{LN}_1(x)),\qquad
  x \leftarrow x + \mathrm{FFN}(\mathrm{LN}_2(x)).
\]
Nhóm chọn pre-LN thay vì post-LN của bài gốc vì gradient đi qua nhánh residual
không bị LayerNorm chắn, nên huấn luyện ổn định hơn ở learning rate cố định mà
không cần warmup. Đây là lựa chọn thiết kế của nhóm, được ghi nhận là khác bài gốc.

FFN theo từng vị trí, hai lớp tuyến tính với ReLU ở giữa
\cite[§3.3]{attention}. Không có cross-attention vì mô hình chỉ có decoder.

\subsubsection{Đối chiếu \dalam}

\begin{center}
\begin{tabular}{p{0.26\textwidth}p{0.30\textwidth}p{0.34\textwidth}}
\toprule
\textbf{Hạng mục} & \textbf{Đối chiếu với} & \textbf{Kết quả}\\
\midrule
A.1a attention & \texttt{F.scaled\_dot\_product\_}\allowbreak\texttt{attention} &
Khớp cả khi có và không có mặt nạ\\
A.1b nhiều đầu & \texttt{nn.MultiheadAttention} cùng trọng số &
Khớp đầu ra \emph{và} gradient theo đầu vào\\
A.1c mặt nạ & Trung bình tiền tố tính tay &
Đạt; kiểm thêm đổi hậu tố và cắt ngắn chuỗi\\
A.1d PE & \texttt{sin}/\texttt{cos} vô hướng &
Khớp bảng đã cộng vào embedding, kể cả $\dmodel$ lẻ\\
A.1e decoder & Chạy tiền tố trên CPU &
Luồng mặt nạ và các thành phần tự cài được dùng trực tiếp\\
\bottomrule
\end{tabular}
\end{center}

Phép kiểm \emph{đổi hậu tố} đáng chú ý: giữ nguyên $y_{\le t}$ và thay đổi
$y_{>t}$, đầu ra tại vị trí $t$ phải không đổi. Phép kiểm này bắt được rò rỉ tương
lai mà phép so sánh hình dạng không bắt được.

%--------------------------------------------------------------------
\subsection{A.2 --- Mô hình ngôn ngữ ký tự và perplexity}
\label{sec:a2}

\subsubsection{Dẫn xuất: hợp lý cực đại trên chuỗi ký tự}

Mục tiêu là cực đại log-hợp lý của dữ liệu dưới \eqref{eq:chain}, tương đương cực
tiểu cross-entropy trung bình theo token. Mục~\ref{sec:b6} chứng minh mục tiêu này
chỉ khác $\KL(\pdata\|p_\theta)$ một hằng số không phụ thuộc $\theta$, nên hai cách
phát biểu cho cùng tập nghiệm. Perplexity là hàm mũ của cross-entropy trung bình
theo token, tức nghịch đảo trung bình nhân của xác suất token đúng.

\subsubsection{Dữ liệu và quyền sử dụng}

Corpus là bản \emph{Truyện Kiều} do Trương Vĩnh Ký phiên âm, in năm 1911, thuộc
phạm vi công cộng theo thời hạn bảo hộ. Nhóm chuyển sang bản này sau khi rà soát
cho thấy căn cứ quyền đối với bản số hóa trước đó (revision \texttt{206607}) chưa
đủ. Toàn bộ quá trình rà soát, kể cả kết luận ``chưa đủ căn cứ'', được giữ trong
\texttt{docs/data/char-corpus.md} thay vì xóa đi.

\begin{center}
\begin{tabular}{lr}
\toprule
\textbf{Chỉ số} & \textbf{Giá trị}\\
\midrule
Số dòng thơ & 3.254\\
Số code point NFC cả corpus & 104.181\\
Tách train (prefix 90\%) & 93.762 code point\\
Tách validation (suffix 10\%) & 10.419 code point\\
Kích thước từ vựng & 127 (gồm UNK, chỉ học từ train)\\
Độ dài ngữ cảnh / stride & 128 / 128\\
Số cửa sổ train / validation & 732 / 81\\
Code point đuôi không dùng (train / validation) & 65 / 50\\
Target được chấm ở validation / số UNK & 10.368 / 1\\
\bottomrule
\end{tabular}
\end{center}

Ba quyết định về dữ liệu đáng nêu, vì mỗi quyết định chặn một đường rò rỉ:

\begin{enumerate}[leftmargin=2em]
\item \textbf{Tách train/validation \emph{trước} khi học từ vựng.} Nếu học từ vựng
trên toàn corpus, thông tin về ký tự chỉ xuất hiện ở validation sẽ lọt vào mô hình
qua bảng embedding. Từ vựng được học chỉ từ train; ký tự validation chưa gặp nhận
ID \texttt{<UNK>} bằng $0$.
\item \textbf{Tách trước khi tạo cửa sổ.} Nếu tạo cửa sổ trước rồi chia, một cửa sổ
có thể bắc qua ranh giới và đưa ký tự validation vào đầu vào train.
\item \textbf{\texttt{<UNK>} vẫn được chấm.} Có đúng một đích \texttt{<UNK>} trong
validation; nó không bị loại khỏi mẫu số perplexity. Loại nó ra sẽ làm PPL tốt hơn
một cách giả tạo.
\end{enumerate}

\subsubsection{Tổng hợp perplexity}

Quy tắc tổng hợp quan trọng ngang với công thức: cộng tổng log-xác suất âm rồi chia
tổng số token được chấm, \emph{không} lấy trung bình các PPL từng batch
(Mục~\ref{sec:b6} giải thích vì sao). Cài đặt tích lũy ở \texttt{float64} với
\texttt{reduction='sum'}:

\begin{quote}\small\ttfamily
nll\_sum += F.cross\_entropy(logits.flatten(0, 1), batch[:, 1:].reshape(-1),\\
\hspace*{6em}reduction='sum')\\
token\_count = windows.size(0) * (windows.size(1) - 1)\\
nll = nll\_sum.item() / token\_count\\
\ldots 'validation\_ppl': math.exp(nll)
\end{quote}

\subsubsection{Kết quả \dalam}

Năm seed độc lập, mỗi seed 5 epoch với 2 decoder block, $\dmodel=64$, 4 đầu,
$d_{ff}=128$, batch size 16, AdamW learning rate $0{,}001$. Các seed dùng cùng một
split; mỗi lượt chấm 10.368 ký tự đích (81 cửa sổ $\times$ 128).

\begin{center}
\begin{tabular}{rrr}
\toprule
\textbf{Seed} & \textbf{NLL validation (nat/ký tự)} & \textbf{PPL validation}\\
\midrule
41 & 2,326985 & 10,2470\\
42 & 2,328789 & 10,2655\\
43 & 2,321075 & 10,1866\\
44 & 2,318854 & 10,1640\\
45 & 2,313612 & 10,1109\\
\midrule
\multicolumn{2}{r}{\textbf{Trung bình $\pm$ SD mẫu}} & $\mathbf{10{,}1948 \pm 0{,}0628}$\\
\bottomrule
\end{tabular}
\end{center}

Con số $\pm 0{,}0628$ là độ lệch chuẩn mẫu (mẫu số $4$) của 5 lần huấn luyện độc
lập \emph{trên cùng một split}. Nó đo độ biến thiên do khởi tạo và thứ tự trộn,
không phải khoảng tin cậy, và không phải đánh giá trên tập test độc lập. Việc này
thỏa yêu cầu tái lập ở phần~II của đề (ít nhất 5 lần chạy, báo cáo trung bình $\pm$
độ lệch chuẩn) cho riêng hạng mục A.2.

\textbf{Giới hạn cần nêu.} Huấn luyện chạy trên Kaggle, phiên hiển thị cấu hình
\texttt{GPU T4 x2}; đây là cấu hình nền tảng hiển thị, không phải xác nhận mô hình
đã dùng cả hai GPU. Tệp \texttt{summary.json} nhận được không chứa phiên bản
PyTorch/CUDA và không kèm checkpoint, nên nhóm \emph{chưa} đối chiếu được log môi
trường. Mục~\ref{sec:tai-lap} ghi việc này là hạng mục còn thiếu.

%--------------------------------------------------------------------
\subsection{A.3 --- Lớp LoRA có hợp nhất trọng số}
\label{sec:a3}

\subsubsection{Dẫn xuất: giới hạn cập nhật vào một không gian hạng thấp}

Tinh chỉnh toàn phần cập nhật cả $W\in\R^{d\times k}$, tức $dk$ tham số cho mỗi
lớp. LoRA ràng buộc cập nhật nằm trong tập các ma trận hạng không quá $r$
\cite[§4.1]{lora}:
\[
  W'=W+\Delta W,\qquad \Delta W=\frac{\alpha}{r}BA,\qquad
  B\in\R^{d\times r},\ A\in\R^{r\times k}.
\]
Đây là bài toán tối ưu trên một đa tạp tham số nhỏ hơn hẳn: thay vì $dk$ chiều,
chỉ còn $r(d+k)$ chiều (Mục~\ref{sec:b4}). Hệ số $\alpha/r$ giữ độ lớn cập nhật
tương đối ổn định khi đổi $r$, nên learning rate không phải điều chỉnh lại theo
từng rank.

\subsubsection{Cài đặt}

Lớp bọc trực tiếp \texttt{nn.Linear} nền và đóng băng \emph{cả weight lẫn bias}:

\begin{quote}\small\ttfamily
self.scale = alpha / rank\\
self.A = nn.Parameter(base.weight.new\_empty(rank, base.in\_features))\\
self.B = nn.Parameter(base.weight.new\_zeros(base.out\_features, rank))\\
nn.init.normal\_(self.A, std=base.in\_features ** -0.5)\\
self.base = base.requires\_grad\_(False)
\end{quote}

Bốn chi tiết có chủ đích:

\begin{enumerate}[leftmargin=2em]
\item $B$ khởi tạo bằng $0$ nên $W'=W$ tại bước đầu, bất kể $A$
(Mục~\ref{sec:b4}). Mô hình xuất phát đúng bằng mô hình nền.
\item $A$ khởi tạo ngẫu nhiên với $\mathrm{std}=k^{-1/2}$. Nếu khởi tạo \emph{cả
hai} bằng $0$ thì gradient theo cả $A$ và $B$ đều triệt tiêu và adapter không bao
giờ học --- câu hỏi thường gặp ở vấn đáp.
\item Lan truyền xuôi \emph{không} tạo ma trận $BA$ kích thước $d\times k$:
\begin{quote}\small\ttfamily
return torch.add(self.base(x), linear(linear(x, self.A), self.B), alpha=self.scale)
\end{quote}
Tính $Ah$ rồi $B(Ah)$ tốn $r(d+k)$ phép nhân, so với $dk$ nếu dựng $BA$ trước.
\item Ràng buộc $1\le r\le\min(d,k)$ được kiểm ở hàm khởi tạo. Chặn trên không cần
cho tính đúng đắn của Mục~\ref{sec:b4}, nhưng $r>\min(d,k)$ không thêm biểu diễn
nào mà vẫn tốn tham số.
\end{enumerate}

Hợp nhất xuất ra một \texttt{nn.Linear} độc lập, giữ bias gốc:

\begin{quote}\small\ttfamily
merged = deepcopy(self.base)\\
merged.weight.addmm\_(self.B, self.A, alpha=self.scale)\\
return merged.eval()
\end{quote}

Hàm trả lớp mới và \emph{không} thay đổi wrapper, nên vẫn học tiếp được sau khi đã
xuất bản suy luận.

\subsubsection{Đối chiếu \dalam}

Notebook A.3 kiểm trên CPU: số tham số học đúng bằng $r(d+k)$; đầu ra của lớp đã
hợp nhất khớp đầu ra hai nhánh \emph{sau khi cập nhật khác $0$} (nếu kiểm lúc
$B$ vẫn bằng $0$ thì phép kiểm vô nghĩa, vì cả hai đều bằng lớp nền); gradient
không chảy vào weight và bias của lớp nền.

%--------------------------------------------------------------------
\subsection{A.4 --- Hàm mất mát DPO}
\label{sec:a4}

\subsubsection{Dẫn xuất}

Mục~\ref{sec:b5} chứng minh hai kết quả rồi ghép lại: nghiệm của bài toán reward
có phạt KL là phân phối Gibbs
$\pi^*\propto\pref\exp(R/\beta)$, và mô hình Bradley--Terry cho xác suất ưu tiên
là $\sigma(R_w-R_l)$. Giải nghịch quan hệ thứ nhất theo reward rồi thế vào quan hệ
thứ hai, hạng $\beta\log Z(x)$ triệt tiêu vì hai đáp án thuộc \emph{cùng} một câu
hỏi, và còn lại
\begin{equation}
  \mathcal L_{\mathrm{DPO}}(\theta)
    =-\E\bigl[\log\sigma(z_\theta)\bigr],\qquad
  z_\theta=\beta\Bigl[
    \log\tfrac{\pi_\theta(y_w\mid x)}{\pref(y_w\mid x)}
    -\log\tfrac{\pi_\theta(y_l\mid x)}{\pref(y_l\mid x)}\Bigr],
  \label{eq:dpo-a4}
\end{equation}
đúng phương trình (7) của \cite{dpo}. Không cần tính $Z(x)$ và không cần huấn
luyện reward model riêng.

\subsubsection{Cài đặt}

\begin{quote}\small\ttfamily
margin = (pw - pl) - (rw - rl)\\
return -logsigmoid(beta * margin).mean()
\end{quote}

Bốn quyết định về tính đúng đắn số học:

\begin{enumerate}[leftmargin=2em]
\item \textbf{Reference không nhận gradient:} \texttt{ref\_chosen.detach()} và
\texttt{ref\_rejected.detach()}. $\pref$ cố định theo định nghĩa bài toán; nếu
gradient chảy vào nó thì \eqref{eq:dpo-a4} không còn là mục tiêu đang tối ưu.
\item \textbf{Dùng \texttt{logsigmoid}, không dùng \texttt{log(sigmoid(\ldots))}.}
Khi $z_\theta$ âm lớn, \texttt{sigmoid} làm tròn xuống $0$ và log trả $-\infty$.
\texttt{logsigmoid} tính ổn định trên toàn miền.
\item \textbf{Không trừ hoặc nhân $\beta$ ở độ chính xác nửa:} tensor \texttt{fp16}
được nâng lên \texttt{float32} trước khi lấy hiệu. Hiệu hai log-xác suất chuỗi là
hiệu hai số lớn gần nhau --- trường hợp xấu nhất của sai số làm tròn.
\item \textbf{$\beta$ phải hữu hạn và dương:} kiểm ở đầu hàm.
Mục~\ref{sec:b5} cần $\beta>0$ ở ba chỗ (để $Z(x)$ hữu hạn, để chia $R_y/\beta$,
và để $\argmax J=\argmin\KL$), nên đây là ràng buộc toán học chứ không phải phòng
vệ đầu vào.
\end{enumerate}

Mỗi đầu vào là \emph{tổng} log-xác suất các token đáp án cho cùng prompt, không
gồm prompt và padding, và \emph{không} chia cho độ dài. Mục~\ref{sec:b6} giải thích
vì sao DPO cộng còn perplexity lại chia.

\subsubsection{Đối chiếu \dalam\ và hạng mục còn thiếu \chualam}

Notebook A.4 so mất mát và \emph{cả hai} gradient theo log-xác suất policy với
cross-entropy hai lớp trên CPU, dùng số minh họa. Bài C.4
(Mục~\ref{sec:c4}) kiểm lại trên ba dấu của margin và bốn giá trị $\beta$, khớp
autograd.

Tuy nhiên yêu cầu G5 của đề --- ít nhất một cài đặt Phần~A được dùng hoặc so sánh
\emph{trực tiếp trong quy trình Phần~D} --- \textbf{chưa} được đáp ứng: hàm này
chưa chạy trên một minibatch dữ liệu ưu tiên thật của Phần~D, vì chưa có dữ liệu
đó. Đây là cầu nối A$\leftrightarrow$D mà nhóm dự định dùng, và nó bị chặn bởi
D.3/D.4 (Mục~\ref{sec:phan-d}).

%--------------------------------------------------------------------
\subsection{Tổng hợp trạng thái Phần A}

\begin{center}
\begin{tabular}{llp{0.42\textwidth}}
\toprule
\textbf{Mã} & \textbf{Trạng thái} & \textbf{Bằng chứng}\\
\midrule
A.1a--A.1e & \dalam & \texttt{src/attention.py}; đối chiếu SDPA và MultiheadAttention\\
A.2 & \dalam & PPL $10{,}1948\pm0{,}0628$, 5 seed; chưa đối chiếu log môi trường\\
A.3 & \dalam & \texttt{src/lora.py}; đếm tham số và kiểm hợp nhất\\
A.4 & \motphan & \texttt{src/dpo.py} đã kiểm số học; chưa chạy trên minibatch thật của D\\
\bottomrule
\end{tabular}
\end{center}

\clearpage
\section{Phần B --- Chứng minh lý thuyết}
\label{sec:phan-b}

Mục này chứng minh sáu mệnh đề của Đề tài 9 \cite[tr.~18]{de}. Mỗi mệnh đề nêu
giả thiết, kích thước các đối tượng, các bước suy luận và kết luận. Các chứng minh
không cần đầu vào của Phần~C. Mọi logarit là logarit tự nhiên. Ví dụ số trong mục
này chỉ để minh họa.

Bản độc lập của mục này nằm ở \texttt{docs/theory/part-b.tex} và biên dịch riêng
được.

%--------------------------------------------------------------------
\subsection{B.1 --- Phương sai tích vô hướng và hệ số chia tỷ lệ}
\label{sec:b1}

\textbf{Mệnh đề.} Cho $d_k\ge 1$ và $q,k\in\R^{d_k}$. Giả sử toàn bộ $2d_k$ thành
phần của $q$ và $k$ độc lập lẫn nhau, mỗi thành phần có kỳ vọng $0$ và phương sai
$1$. Khi đó $\Var(q^\top k)=d_k$. Hệ số $1/\sqrt{d_k}$ giữ phương sai của mỗi điểm
chú ý đã chia tỷ lệ bằng $1$ trong mô hình giả thiết này
\cite[§3.2.1, chú thích 4]{attention}.

\subsubsection{Tính phương sai}

Đặt $Z_i=q_i k_i$ với $1\le i\le d_k$. Từ $\Var(q_i)=\E[q_i^2]-(\E[q_i])^2$ và
$\E[q_i]=0$, suy ra $\E[q_i^2]=1$; tương tự $\E[k_i^2]=1$. Vì $q_i$ và $k_i$ độc
lập,
\[
  \E[Z_i]=\E[q_i]\E[k_i]=0,\qquad
  \E[Z_i^2]=\E[q_i^2]\E[k_i^2]=1,
\]
nên $\Var(Z_i)=1$. Với $i\ne j$, hai biến $Z_i$ và $Z_j$ là hàm của hai nhóm biến
độc lập $(q_i,k_i)$ và $(q_j,k_j)$, do đó $\Cov(Z_i,Z_j)=0$. Vì
$q^\top k=\sum_{i=1}^{d_k} Z_i$,
\begin{align}
  \E[q^\top k]&=0,\notag\\
  \Var(q^\top k)
    &=\sum_{i=1}^{d_k}\Var(Z_i)+2\!\!\sum_{1\le i<j\le d_k}\!\!\Cov(Z_i,Z_j)=d_k,
    \label{eq:b1-var}\\
  \Var\!\left(\frac{q^\top k}{\sqrt{d_k}}\right)
    &=\frac{1}{d_k}\Var(q^\top k)=1.\notag
\end{align}

Chú ý rằng giả thiết độc lập được dùng \emph{hai lần} với hai vai trò khác nhau:
lần thứ nhất để tách $\E[q_ik_i]$ thành tích hai kỳ vọng, lần thứ hai để triệt tiêu
hiệp phương sai giữa các chỉ số khác nhau. Chỉ giả thiết ``không tương quan'' là
không đủ cho lần thứ nhất.

\subsubsection{Vì sao điều này giúp giảm bão hòa softmax?}

Xét $m\ge2$ điểm chú ý $z_1,\ldots,z_m$ trong một hàng, với
$s_i=e^{z_i}/\sum_{j=1}^m e^{z_j}$. Dùng quy tắc đạo hàm thương, với
$\delta_{ij}=1$ nếu $i=j$ và $0$ nếu $i\ne j$:
\begin{equation}
  \frac{\partial s_i}{\partial z_j}
  =\frac{\delta_{ij}e^{z_i}}{\sum_a e^{z_a}}
   -\frac{e^{z_i}e^{z_j}}{(\sum_a e^{z_a})^2}
  =s_i(\delta_{ij}-s_j).
  \label{eq:b1-jac}
\end{equation}
Nếu điểm lớn nhất $z_a$ hơn mọi điểm khác ít nhất $g>0$ thì
\begin{equation}
  s_a=\frac{1}{1+\sum_{j\ne a} e^{-(z_a-z_j)}}
    \ \ge\ \frac{1}{1+(m-1)e^{-g}}.
  \label{eq:b1-sat}
\end{equation}
Giữ $m$ cố định và cho $g$ tăng, $s_a\to1$, các $s_j$ còn lại $\to0$, và theo
\eqref{eq:b1-jac} toàn bộ Jacobian của softmax tiến đến $0$. Khi đó trọng số chú ý
gần như không phản ứng với nhiễu nhỏ ở điểm số, tức gradient qua lớp attention bị
triệt tiêu.

Điểm then chốt: bão hòa phụ thuộc vào \emph{chênh lệch} điểm, không phụ thuộc việc
cộng một hằng số chung vào cả hàng --- softmax bất biến với phép dịch đó. Với các
tích vô hướng thô, độ lệch chuẩn của từng điểm là $\sqrt{d_k}$ theo
\eqref{eq:b1-var}. Chia cả hàng cho $\sqrt{d_k}$ đưa độ lệch chuẩn mỗi điểm về $1$
và thu nhỏ mọi chênh lệch theo cùng hệ số khi $d_k>1$.

\textbf{Phạm vi kết luận.} Kết quả phương sai giải thích \emph{lý do chọn} thang
điểm. Nó không khẳng định mọi hàng đều không bão hòa, cũng không khẳng định các tọa
độ của query và key trong mô hình đã huấn luyện là độc lập hay có phương sai đơn
vị. Giả thiết ở đây là một mô hình giả định tại khởi tạo.

%--------------------------------------------------------------------
\subsection{B.2 --- Tính đẳng biến hoán vị của self-attention}
\label{sec:b2}

\textbf{Mệnh đề.} Xét self-attention không mã hóa vị trí, không mặt nạ, dùng cùng
trọng số chiếu cho mọi hàng. Với $X\in\R^{n\times \dmodel}$ mà mỗi hàng là biểu
diễn một token, nếu $P\in\R^{n\times n}$ là ma trận hoán vị thì
$\Attn(PX)=P\Attn(X)$.

Đây là tính \emph{đẳng biến} (equivariance): đầu ra đổi thứ tự theo đầu vào. Nó
khác tính \emph{bất biến} (invariance), trong đó đầu ra không đổi.

\subsubsection{Theo dõi phép hoán vị đến ma trận điểm}

Dùng $W_Q,W_K\in\R^{\dmodel\times d_k}$ và $W_V\in\R^{\dmodel\times d_v}$, với
$n,\dmodel,d_k,d_v>0$. Đặt
\[
  Q=XW_Q,\quad K=XW_K,\quad V=XW_V,\qquad
  S=\frac{QK^\top}{\sqrt{d_k}}\in\R^{n\times n}.
\]
Theo \eqref{eq:sdpa}, $\Attn(X)=\softmax_{\mathrm{row}}(S)V$. Với đầu vào $PX$, vì
các phép chiếu tác động \emph{bên phải} còn hoán vị tác động \emph{bên trái}, ta có
$Q'=PQ$, $K'=PK$, $V'=PV$, và do đó
\[
  S'=\frac{(PQ)(PK)^\top}{\sqrt{d_k}}=P\,\frac{QK^\top}{\sqrt{d_k}}\,P^\top=PSP^\top.
\]

\subsubsection{Hoán vị đi qua softmax theo hàng}

Chọn quy ước $(PX)_{i,:}=X_{\tau(i),:}$ với $\tau$ là song ánh trên
$\{1,\ldots,n\}$. Khi đó $(PSP^\top)_{ij}=S_{\tau(i),\tau(j)}$ và
\begin{align}
  [\softmax_{\mathrm{row}}(PSP^\top)]_{ij}
    &=\frac{\exp(S_{\tau(i),\tau(j)})}{\sum_{a=1}^{n}\exp(S_{\tau(i),\tau(a)})}
     \notag\\
    &=\frac{\exp(S_{\tau(i),\tau(j)})}{\sum_{b=1}^{n}\exp(S_{\tau(i),b})}
     \label{eq:b2-key}\\
    &=[P\softmax_{\mathrm{row}}(S)P^\top]_{ij}.\notag
\end{align}
Đẳng thức \eqref{eq:b2-key} là bước quyết định: nó đúng vì $\tau$ là song ánh, nên
$a\mapsto\tau(a)$ đi qua đúng một lần mỗi chỉ số cột và tổng ở mẫu số không đổi
giá trị. Nếu $\tau$ không phải song ánh, bước này sai.

Do $P^\top P=I_n$,
\begin{align*}
  \Attn(PX)
    &=\softmax_{\mathrm{row}}(PSP^\top)\,PV
     =P\softmax_{\mathrm{row}}(S)P^\top PV\\
    &=P\softmax_{\mathrm{row}}(S)V=P\Attn(X).
\end{align*}

Trong chú ý nhiều đầu, lập luận trên đúng cho từng đầu. Ghép các đầu theo chiều đặc
trưng rồi nhân bên phải với $W_O$ vẫn cho đúng phép hoán vị hàng, vì ghép các
$PH_i$ bằng $P$ nhân với ma trận ghép các $H_i$.

\subsubsection{Giới hạn: mặt nạ nhân quả phá đẳng biến}

Với mặt nạ cộng $M\in(\R\cup\{-\infty\})^{n\times n}$ có ít nhất một vị trí không
bị che ở mỗi hàng, đặt $\Attn_M(X)=\softmax_{\mathrm{row}}(S+M)V$. Lập luận trên
vẫn cho
\begin{equation}
  \Attn_{PMP^\top}(PX)=P\Attn_M(X),
  \label{eq:b2-mask}
\end{equation}
nghĩa là \emph{nếu hoán vị cả mặt nạ} thì đẳng biến được giữ.

Điều này \emph{không} cho phép giữ nguyên một mặt nạ nhân quả cố định. Phản ví dụ:
lấy hai token vô hướng, $X=(0,2)^\top$, $W_Q=W_K=0$, $W_V=1$, và $P$ đổi chỗ hai
hàng. Vì $W_Q=W_K=0$, mọi điểm số bằng $0$, nên softmax chia đều trên các khóa
được phép. Giữ mặt nạ nhân quả (hàng một nhìn chính nó; hàng hai nhìn cả hai):
\[
  \Attn_M(X)=(0,1)^\top,\qquad
  \Attn_M(PX)=(2,1)^\top\ \ne\ (1,0)^\top=P\Attn_M(X).
\]
Vậy mệnh đề phổ quát ``attention có mặt nạ nhân quả cố định là đẳng biến'' bị bác
bỏ. Nguyên nhân cấu trúc: mặt nạ nhân quả gắn với \emph{vị trí}, không gắn với
token, nên đổi thứ tự token mà giữ mặt nạ là đổi bài toán.

Tương tự, khi cộng bảng mã hóa vị trí cố định $E$, đổi thứ tự token nhưng giữ
nguyên vị trí cho $PX+E$, nói chung khác $P(X+E)$. Vì thế định lý về attention
không PE không suy ra đẳng biến cho mô hình thực tế có PE và có mặt nạ.

\textbf{Về phản ví dụ suy biến.} Phản ví dụ dùng $d_{\mathrm{model}}=1$ và phép
chiếu bằng $0$. Điều này hợp lệ: để bác bỏ một mệnh đề phổ quát chỉ cần một ví dụ,
và kết quả dương \eqref{eq:b2-mask} vẫn được giữ nguyên bên cạnh.

%--------------------------------------------------------------------
\subsection{B.3 --- Ma trận quay khối của mã hóa vị trí hình sin}
\label{sec:b3}

\textbf{Mệnh đề.} Xét vector cột $\PE(p)\in\R^{2m}$ gồm $m\ge1$ cặp đầy đủ theo
thứ tự sin trước, cos sau: với tần số cố định $\omega_i$, $i=0,\ldots,m-1$, đặt
$\PE(p)_{2i}=\sin(p\omega_i)$ và $\PE(p)_{2i+1}=\cos(p\omega_i)$. Khi đó tồn tại
$M_\Delta\in\R^{2m\times 2m}$ chỉ phụ thuộc độ dịch $\Delta$ (với hệ tần số đã cố
định) sao cho $\PE(p+\Delta)=M_\Delta\PE(p)$ với mọi $p$. Trong Transformer gốc,
$\omega_i=10000^{-2i/\dmodel}$ \cite[§3.5]{attention}.

\subsubsection{Chứng minh cho một cặp tọa độ}

Đặt $a=p\omega_i$ và $b=\Delta\omega_i$. Công thức cộng góc cho
\begin{align*}
  \sin(a+b)&=\cos b\,\sin a+\sin b\,\cos a,\\
  \cos(a+b)&=-\sin b\,\sin a+\cos b\,\cos a.
\end{align*}
Vậy nếu $u_i(p)=(\sin(p\omega_i),\cos(p\omega_i))^\top$ thì
$u_i(p+\Delta)=M_{\Delta,i}\,u_i(p)$ với
\begin{equation}
  M_{\Delta,i}=
    \begin{pmatrix}
      \cos(\Delta\omega_i)&\sin(\Delta\omega_i)\\
      -\sin(\Delta\omega_i)&\cos(\Delta\omega_i)
    \end{pmatrix}.
  \label{eq:b3-block}
\end{equation}
Hai cột của khối này có chuẩn $1$ vì $\cos^2 b+\sin^2 b=1$; tích vô hướng của chúng
bằng $\cos b\sin b-\sin b\cos b=0$; định thức bằng $\cos^2 b+\sin^2 b=1$. Do đó
\eqref{eq:b3-block} là ma trận quay --- theo quy ước vector cột $(\sin,\cos)^\top$,
nó là phép quay góc $-b$ trong mặt phẳng tọa độ chuẩn.

\subsubsection{Ghép các cặp thành vector đầy đủ}

Lấy $M_\Delta=\diag(M_{\Delta,0},\ldots,M_{\Delta,m-1})$. Vì các khối ngoài đường
chéo bằng $0$, phép nhân tác động riêng lên từng cặp tọa độ, đúng bằng
$u_i(p)\mapsto u_i(p+\Delta)$. Ghép lại theo thứ tự ban đầu được
$M_\Delta\PE(p)=\PE(p+\Delta)$. Các phần tử của $M_\Delta$ chỉ chứa $\Delta$ và
các tần số đã cố định, \emph{không} chứa $p$ --- đây là điều cần chứng minh.

Kết quả cho thấy độ dịch cố định có \emph{biểu diễn tuyến tính} trên PE. Nó không
khẳng định mô hình đã học được cách khai thác biểu diễn đó.

\subsubsection{Phạm vi: mệnh đề cần số chiều chẵn}
\label{sec:b3-odd}

Giả thiết ``$m$ cặp đầy đủ'' là điều kiện thực sự cần, không phải quy ước trình
bày. Cài đặt \texttt{SinusoidalPositionEncoding} của nhóm chấp nhận $\dmodel$ lẻ:
với $\dmodel=2m+1$, mã sinh $m+1$ tần số, đặt $\sin$ vào $m+1$ cột chỉ số chẵn và
$\cos$ vào $m$ cột chỉ số lẻ, nên
\[
  \PE(p)=\bigl(\sin p\omega_0,\ \cos p\omega_0,\ \ldots,\
                \sin p\omega_{m-1},\ \cos p\omega_{m-1},\
                \sin p\omega_m\bigr)^\top .
\]
Tọa độ cuối không có $\cos(p\omega_m)$ đi kèm. Ta chứng minh khi đó \emph{không tồn
tại} $M_\Delta$ nào thỏa mệnh đề, trừ khi $\sin(\Delta\omega_m)=0$.

\textbf{Bổ đề.} Cho các tần số phân biệt $\omega_0,\ldots,\omega_m\in(0,\pi)$. Nếu
$L\ge 2m+2$ thì $2m+2$ hàm $p\mapsto\sin(p\omega_i)$ và $p\mapsto\cos(p\omega_i)$,
$0\le i\le m$, độc lập tuyến tính trên $\R$ khi xét như hàm trên
$\{0,1,\ldots,L-1\}$.

\emph{Chứng minh.} Đặt $z_i=e^{\mathrm i\omega_i}$. Mỗi tổ hợp tuyến tính thực của
các hàm trên viết được thành
$\sum_{i=0}^{m}\bigl(\gamma_i z_i^{\,p}+\overline{\gamma_i}\,\overline{z_i}^{\,p}\bigr)$.
Vì các $\omega_i\in(0,\pi)$ phân biệt, $2m+2$ số
$z_0,\overline{z_0},\ldots,z_m,\overline{z_m}$ đôi một khác nhau. Nếu tổng triệt
tiêu tại $p=0,\ldots,2m+1$ thì hệ Vandermonde với các nút phân biệt đó khả nghịch,
buộc mọi $\gamma_i=0$. \hfill$\square$

Với lựa chọn của Transformer gốc, $\omega_i=10000^{-2i/\dmodel}\in(0,1]\subset(0,\pi)$
và đôi một phân biệt, nên bổ đề áp dụng khi $\texttt{max\_len}\ge\dmodel+1$.

\textbf{Suy ra.} Giả sử có $M_\Delta\in\R^{(2m+1)\times(2m+1)}$ với
$\PE(p+\Delta)=M_\Delta\PE(p)$ cho mọi $p$ trong miền vị trí. Đọc hàng cuối (chỉ số
$2m$) và dùng công thức cộng góc:
\[
  \cos(\Delta\omega_m)\sin(p\omega_m)+\sin(\Delta\omega_m)\cos(p\omega_m)
  =\sum_{i=0}^{m-1}\bigl(a_i\sin(p\omega_i)+b_i\cos(p\omega_i)\bigr)
   +a_m\sin(p\omega_m),
\]
với $a_i,b_i$ là các phần tử của hàng đó, không phụ thuộc $p$. Chuyển vế:
\[
  \sin(\Delta\omega_m)\cos(p\omega_m)
  =\sum_{i=0}^{m-1}\bigl(a_i\sin(p\omega_i)+b_i\cos(p\omega_i)\bigr)
   +\bigl(a_m-\cos(\Delta\omega_m)\bigr)\sin(p\omega_m).
\]
Theo bổ đề, $\cos(p\omega_m)$ độc lập tuyến tính với mọi hàm ở vế phải, nên đẳng
thức chỉ đúng với mọi $p$ khi $\sin(\Delta\omega_m)=0$, tức
$\Delta\omega_m\in\pi\mathbb Z$. Vì $0<\omega_m\le1$, điều này loại mọi độ dịch
nguyên $\Delta$ với $|\Delta|<\pi/\omega_m$, trong đó có $\Delta=1$. Mệnh đề thất
bại ở số chiều lẻ.

\textbf{Hệ quả cho cài đặt.} Thí nghiệm A.2 dùng $\dmodel=64$ (chẵn) nên bảng PE
được dùng thực tế vẫn thỏa mệnh đề; phần kiểm thử $\dmodel$ lẻ của nhóm chỉ đối
chiếu giá trị bảng, không kiểm tính chất dịch vị trí. Muốn giữ tính chất cho số
chiều lẻ, cần bỏ tọa độ $\sin$ lẻ cặp hoặc đệm lên số chiều chẵn; đó là lựa chọn
thiết kế, không phải hệ quả của chứng minh.

%--------------------------------------------------------------------
\subsection{B.4 --- Khởi tạo, hạng, số tham số và hợp nhất LoRA}
\label{sec:b4}

\textbf{Thiết lập.} Lớp tuyến tính có $W\in\R^{d\times k}$ với $d,k\ge1$, đầu vào
cột $h\in\R^k$. LoRA đóng băng $W$, học $A\in\R^{r\times k}$ và $B\in\R^{d\times r}$
với $r\ge1$ nguyên và $\alpha$ cố định. Không huấn luyện bias; nếu có bias cố định
$b\in\R^d$, nó được cộng như nhau trong cả hai cách tính. Đặt
$\Delta W=(\alpha/r)BA$ và $W'=W+\Delta W$ \cite[§4.1--4.2]{lora}.

\subsubsection{Khởi tạo không thay đổi lớp gốc}

Khi $B=0_{d\times r}$, mỗi phần tử $(BA)_{ij}=\sum_{a=1}^{r}B_{ia}A_{aj}$ bằng $0$
bất kể giá trị của $A$. Do đó $\Delta W=0_{d\times k}$ và $W'=W$. Bài báo LoRA chọn
$A$ ngẫu nhiên và $B=0$; kết luận $W'=W$ chỉ cần điều kiện thứ hai.

Ngược lại, nếu khởi tạo \emph{cả hai} bằng $0$ thì
$\partial\mathcal L/\partial A\propto B^\top(\cdot)=0$ và
$\partial\mathcal L/\partial B\propto(\cdot)(Ah)^\top=0$, nên adapter không bao giờ
rời khỏi điểm $0$. Vì vậy một trong hai ma trận phải khác $0$ lúc khởi tạo.

\subsubsection{Hạng của cập nhật bị chặn bởi $r$}

Xem $A$ là ánh xạ $\R^k\to\R^r$ và $B$ là ánh xạ $\R^r\to\R^d$. Với mọi $h$,
$BAh=B(Ah)\in\im(B)$. Vì $\im(B)$ sinh bởi $r$ cột của $B$,
\[
  \rank(BA)=\dim\im(BA)\le\dim\im(B)=\rank(B)\le r,
\]
suy ra $\rank(\Delta W)\le r$. Mệnh đề chặn hạng của \emph{phần cập nhật}, không
chặn hạng của $W'$: tổng của một ma trận hạng đầy và một ma trận hạng thấp nói
chung vẫn hạng đầy.

\subsubsection{Đếm tham số thực sự được huấn luyện}

Học trực tiếp toàn bộ $W$ cần cập nhật $dk$ hệ số. Với LoRA, $A$ có $rk$ phần tử và
$B$ có $dr$ phần tử; vì $W$, $\alpha$ và bias không được huấn luyện, số tham số cập
nhật là
\begin{equation}
  rk+dr=r(d+k).
  \label{eq:b4-count}
\end{equation}
Số này nhỏ hơn $dk$ khi $r<dk/(d+k)$; tỷ lệ chính xác so với học cả $W$ là
$r(d+k)/(dk)$. Đây không phải tổng dung lượng mô hình, vì vẫn cần lưu trọng số gốc.

\subsubsection{Hợp nhất không thêm phép tính trong lan truyền xuôi}

Trong suy luận xác định với adapter đã cố định, hai nhánh cho cùng ánh xạ như một
lớp tuyến tính duy nhất:
\begin{equation}
  Wh+\frac{\alpha}{r}B(Ah)=\left(W+\frac{\alpha}{r}BA\right)h=W'h.
  \label{eq:b4-merge}
\end{equation}
Đây là tính phân phối của phép nhân ma trận, không phải xấp xỉ. Tính và lưu $W'$
một lần trước khi phục vụ; sau đó mỗi đầu vào dùng đúng một phép nhân
ma trận--vector $d\times k$ như lớp gốc ($dk$ phép nhân và $d(k-1)$ phép cộng nếu
tính trực tiếp, cộng cùng chi phí bias nếu có). Không còn các phép nhân phụ $Ah$,
$B(Ah)$ hay phép cộng hai nhánh trên mỗi đầu vào.

\textbf{Phạm vi kết luận.} Kết luận nói về chi phí số học \emph{sau hợp nhất}, với
giả thiết cùng kiểu thực thi lớp tuyến tính. Nó không bao gồm chi phí tạo $W'$,
không bảo đảm khớp từng bit trong số học dấu phẩy động, và không tự mở rộng thành
kết luận về độ trễ khi đổi kernel hay tái lượng tử hóa.

\subsubsection{Đối chiếu với cài đặt A.3}

Bốn tiểu mệnh đề trên khớp từng dòng với \texttt{src/lora.py}: $B$ khởi tạo bằng
\texttt{new\_zeros}; hệ số \texttt{scale = alpha / rank}; lớp nền gọi
\texttt{requires\_grad\_(False)} nên cả weight và bias bị đóng băng, đúng giả thiết
đếm tham số \eqref{eq:b4-count}; và \texttt{merge()} dùng
\texttt{addmm\_(B, A, alpha=scale)} trên bản sao, cho đúng
$W+(\alpha/r)BA$ với bias gốc, như \eqref{eq:b4-merge}. Cài đặt ràng buộc thêm
$r\le\min(d,k)$, điều mà chứng minh không cần.

%--------------------------------------------------------------------
\subsection{B.5 --- Nghiệm tối ưu có phạt KL và hàm mất mát DPO}
\label{sec:b5}

\textbf{Thiết lập.} Cố định câu hỏi $x$. Ký hiệu $R(x,y)$ là reward của đáp án $y$
(đề dùng $r(x,y)$), $\beta>0$, và $\pref(\cdot\mid x)$ là phân phối tham chiếu cố
định. Xét tập đáp án hữu hạn, không rỗng $\mathcal Y_x$, với reward hữu hạn trên
hỗ trợ $S_x=\{y\in\mathcal Y_x:\pref(y\mid x)>0\}$. Tối ưu trên mọi phân phối
$\pi(\cdot\mid x)$ có hỗ trợ nằm trong $S_x$; chính sách gán khối lượng dương cho
$y\notin S_x$ có KL vô hạn nên không thể là nghiệm hữu hạn. Quy ước: hạng tử có hệ
số xác suất bằng $0$ trong KL bằng $0$. Theo
\cite[§3--4, phương trình (1)--(7), Phụ lục A.1--A.2]{dpo}.

\subsubsection{Xây dựng một phân phối ứng viên}

Viết tắt $\pi_y=\pi(y\mid x)$, $q_y=\pref(y\mid x)$, $R_y=R(x,y)$ trên $S_x$. Hàm
mục tiêu là
\begin{equation}
  J_x(\pi)=\sum_{y\in S_x}\pi_y R_y
    -\beta\!\!\sum_{y\in S_x:\,\pi_y>0}\!\!\pi_y\log\frac{\pi_y}{q_y}.
  \label{eq:b5-obj}
\end{equation}
Đặt $Z(x)=\sum_{y\in S_x}q_y\exp(R_y/\beta)$. Vì tập hữu hạn, hỗ trợ không rỗng và
các số hạng dương hữu hạn, $0<Z(x)<\infty$. Định nghĩa
\begin{equation}
  \pi^*(y\mid x)=\frac{\pref(y\mid x)\exp(R(x,y)/\beta)}{Z(x)}
  \quad(y\in S_x),
  \label{eq:b5-gibbs}
\end{equation}
và $\pi^*(y\mid x)=0$ ngoài $S_x$. Tổng tử số trên $S_x$ chính bằng $Z(x)$, nên
$\pi^*$ là một phân phối, dương trên toàn bộ $S_x$.

\subsubsection{Chứng minh tối ưu toàn cục và tính duy nhất}

Ký hiệu $\pi_y^*=\pi^*(y\mid x)$. Từ \eqref{eq:b5-gibbs},
$\log(\pi_y^*/q_y)=R_y/\beta-\log Z(x)$. Thế vào và dùng $\sum_y\pi_y=1$:
\begin{align}
  \beta\KL(\pi\|\pi^*)
     &=\beta\!\!\sum_{y:\,\pi_y>0}\!\!\pi_y\log\frac{\pi_y}{\pi_y^*}\notag\\
     &=\beta\!\!\sum_{y:\,\pi_y>0}\!\!\pi_y\log\frac{\pi_y}{q_y}
       -\sum_y\pi_y R_y+\beta\log Z(x)\notag\\
     &=\beta\log Z(x)-J_x(\pi).
  \label{eq:b5-central}
\end{align}
Đẳng thức \eqref{eq:b5-central} là trục của chứng minh: cực đại hóa $J_x$ tương
đương cực tiểu hóa $\KL(\pi\|\pi^*)$, vì $\beta\log Z(x)$ không phụ thuộc $\pi$.

Để thấy điều kiện dấu bằng mà không cần KKT, dùng $-\log t\ge 1-t$ với $t>0$, dấu
bằng chỉ tại $t=1$. Với $T=\{y:\pi_y>0\}$, áp dụng cho $t=\pi_y^*/\pi_y$:
\begin{equation}
  \KL(\pi\|\pi^*)\ \ge\ \sum_{y\in T}(\pi_y-\pi_y^*)
       \ =\ 1-\sum_{y\in T}\pi_y^*\ \ge\ 0.
  \label{eq:b5-unique}
\end{equation}
Nếu $\KL=0$, bất đẳng thức đầu buộc $\pi_y^*=\pi_y$ trên $T$, và bất đẳng thức sau
buộc $\pi^*$ không gán khối lượng ngoài $T$. Ngược lại, $\pi=\pi^*$ cho $\KL=0$.
Vì $\beta>0$, hàm $J_x$ đạt cực đại $\beta\log Z(x)$ \emph{đúng tại} $\pi^*$; không
còn nghiệm nào khác, kể cả ở biên của miền phân phối.

\textbf{Miền đếm được.} Một điều kiện đủ để dùng nguyên lập luận là reward bị chặn
hai phía trên hỗ trợ: nếu $a\le R_y\le b$ thì
$0<e^{a/\beta}\le Z(x)\le e^{b/\beta}<\infty$, nên kỳ vọng reward hữu hạn và
$\log(\pi_y^*/q_y)$ bị chặn. Đẳng thức \eqref{eq:b5-central} vẫn đúng khi KL nhận
giá trị $+\infty$, cho mục tiêu $-\infty$ chứ không tạo hiệu $\infty-\infty$. Nếu
reward không bị chặn, phải kiểm riêng việc chuẩn hóa và tính xác định của mục tiêu.

\subsubsection{Từ Bradley--Terry tới xác suất ưu tiên}

Xét cặp $y_w,y_l\in S_x$ cho cùng $x$, trong đó $y_w$ được ưu tiên. Mô hình
Bradley--Terry gán
\begin{align}
  \Pr(y_w\succ y_l\mid x,y_w,y_l)
   &=\frac{\exp(R(x,y_w))}{\exp(R(x,y_w))+\exp(R(x,y_l))}\notag\\
   &=\sigma\bigl(R(x,y_w)-R(x,y_l)\bigr),
   \qquad \sigma(z)=\frac{1}{1+e^{-z}}.
  \label{eq:b5-bt}
\end{align}
Giải \eqref{eq:b5-gibbs} theo reward:
\begin{equation}
  R(x,y)=\beta\log\frac{\pi^*(y\mid x)}{\pref(y\mid x)}+\beta\log Z(x).
  \label{eq:b5-invert}
\end{equation}
Lấy hiệu giữa hai đáp án \emph{cùng} $x$, hai hạng $\beta\log Z(x)$ triệt tiêu. Do
đó xác suất ưu tiên chỉ còn phụ thuộc hiệu hai log-tỷ số chính sách tối ưu trên
tham chiếu. Bước triệt tiêu \emph{không} áp dụng cho hai đáp án của hai câu hỏi
khác nhau --- đây là lý do dữ liệu ưu tiên phải ghép theo cùng một prompt.

\subsubsection{Hợp lý cực đại theo chính sách tham số hóa}

Dùng $\pi_\theta$ thay cho $\pi^*$ chưa biết trong \eqref{eq:b5-bt}. Hai đáp án của
mỗi cặp phải có xác suất dương dưới \emph{cả} $\pi_\theta$ và $\pref$ để log-tỷ số
hữu hạn. Đặt
\begin{equation}
  z_\theta(x,y_w,y_l)
   =\beta\left[
      \log\frac{\pi_\theta(y_w\mid x)}{\pref(y_w\mid x)}
      -\log\frac{\pi_\theta(y_l\mid x)}{\pref(y_l\mid x)}
   \right].
  \label{eq:b5-logit}
\end{equation}
Với tập $\mathcal D$ các bộ ba đã sắp theo nhãn ưu tiên, mỗi quan sát đóng góp xác
suất $\sigma(z_\theta)$ vào hợp lý của nhãn, có điều kiện trên câu hỏi và cặp đáp
án. Với các quan sát độc lập có điều kiện, hợp lý là tích các xác suất đó; lấy log,
đổi dấu và lấy trung bình trên $\mathcal D$:
\begin{equation}
  \boxed{\ \mathcal L_{\mathrm{DPO}}(\theta)
    =-\E_{(x,y_w,y_l)\sim\mathcal D}\bigl[\log\sigma(z_\theta(x,y_w,y_l))\bigr].\ }
  \label{eq:b5-loss}
\end{equation}

Không cần tính $Z(x)$ và không cần huấn luyện reward model riêng để tính mất mát.
Kiểm dấu: nếu hai log-tỷ số trong \eqref{eq:b5-logit} bằng nhau thì mất mát một cặp
bằng $-\log(1/2)=\log 2\approx0{,}693$; nếu hiệu tăng thì $\sigma$ tăng và mất mát
giảm.

\textbf{Phân biệt hai loại duy nhất.} Tính duy nhất ở \eqref{eq:b5-unique} là duy
nhất của $\pi^*$ \emph{với reward đã cho} khi tối ưu trên các phân phối. Nó không
phải tính duy nhất của tham số $\theta$ khi học từ hữu hạn cặp.

\subsubsection{Đối chiếu với cài đặt A.4}

Biểu thức \eqref{eq:b5-logit} khớp đúng \texttt{src/dpo.py}: mã tính
\texttt{margin = (pw - pl) - (rw - rl)} rồi \texttt{-logsigmoid(beta * margin)
.mean()}, tức $-\log\sigma(z_\theta)$ lấy trung bình trên các cặp, đúng
\eqref{eq:b5-loss}. Ràng buộc $\beta$ hữu hạn và dương trong mã tương ứng với giả
thiết $\beta>0$ ở đây.

%--------------------------------------------------------------------
\subsection{B.6 --- Cross-entropy, KL và perplexity}
\label{sec:b6}

\textbf{Thiết lập.} Phân phối dữ liệu cố định $\pdata$ và họ phân phối mô hình
$p_\theta$ trên cùng miền rời rạc hữu hạn $\mathcal Y$. Giả sử $p_\theta(y)>0$ với
mọi $y$ sao cho $\pdata(y)>0$ khi cần giá trị hữu hạn. Hạng tử có $\pdata(y)=0$
được bỏ khỏi tổng; nếu $\pdata(y)>0$ nhưng $p_\theta(y)=0$ thì cả cross-entropy và
KL đều là $+\infty$.

\subsubsection{Phân rã cross-entropy kỳ vọng}

Đặt $S=\{y:\pdata(y)>0\}$. Theo định nghĩa,
\begin{align*}
  H(\pdata,p_\theta)&=-\sum_{y\in S}\pdata(y)\log p_\theta(y),\\
  H(\pdata)&=-\sum_{y\in S}\pdata(y)\log\pdata(y).
\end{align*}
Khai triển tỷ số trong KL thành hiệu hai logarit:
\begin{equation}
  \KL(\pdata\|p_\theta)
   =\sum_{y\in S}\pdata(y)\log\frac{\pdata(y)}{p_\theta(y)}
   =-H(\pdata)+H(\pdata,p_\theta).
  \label{eq:b6-decomp}
\end{equation}
Vì $H(\pdata)$ hữu hạn và không phụ thuộc $\theta$, hai mục tiêu khác nhau đúng một
hằng số: chúng xếp thứ tự các mô hình như nhau và có cùng tập nghiệm cực tiểu (nếu
cực tiểu đạt được):
\begin{equation}
  \argmin_\theta H(\pdata,p_\theta)=\argmin_\theta \KL(\pdata\|p_\theta).
  \label{eq:b6-equiv}
\end{equation}
Trên miền đếm được, lập luận trừ hằng số vẫn dùng được khi $H(\pdata)$ hữu hạn.
Không được lấy hiệu $\infty-\infty$ nếu entropy dữ liệu vô hạn.

\subsubsection{Phiên bản có ngữ cảnh cho dự đoán token}

Ký hiệu $C$ là ngữ cảnh và $Y$ là token đích. Giữ phân phối của $C$ cố định, áp
dụng \eqref{eq:b6-decomp} cho mỗi ngữ cảnh $c$ rồi lấy kỳ vọng:
\begin{align}
  \E_{(C,Y)\sim\pdata}[-\log p_\theta(Y\mid C)]
   &= H_{\pdata}(Y\mid C)\notag\\
   &\quad+\E_{C\sim\pdata}\bigl[\KL(\pdata(\cdot\mid C)\|p_\theta(\cdot\mid C))\bigr],
  \label{eq:b6-ctx}
\end{align}
trong đó $H_{\pdata}(Y\mid C)=\E_{C}[H(\pdata(\cdot\mid C))]$ không phụ thuộc mô
hình và hữu hạn vì từ vựng hữu hạn. Đây là phát biểu tương ứng khi mô hình ngôn ngữ
được huấn luyện dự đoán token tiếp theo theo ngữ cảnh đúng từ dữ liệu.

\subsubsection{Từ log-xác suất trung bình theo token đến PPL}

Trên tập đánh giá, gọi $(c_t,y_t)$ là ngữ cảnh và token đích ở vị trí được chấm thứ
$t$, tổng cộng $N>0$ token. Với logarit tự nhiên và xác suất dương tại các vị trí
này,
\begin{equation}
  \bar H=-\frac{1}{N}\sum_{t=1}^{N}\log p_\theta(y_t\mid c_t),\qquad
  \PPL=e^{\bar H}
    =\left(\prod_{t=1}^{N} p_\theta(y_t\mid c_t)\right)^{-1/N}.
  \label{eq:b6-ppl}
\end{equation}
Đẳng thức cuối dùng $\exp(\sum_t\log a_t)=\prod_t a_t$. Vậy PPL là nghịch đảo
trung bình nhân của xác suất token đúng; vì hàm mũ tăng chặt, cross-entropy theo
token thấp hơn tương đương PPL thấp hơn \emph{trên cùng thiết lập chấm}.

\emph{Ví dụ minh họa.} Hai token có xác suất đúng $1/2$ và $1/8$ cho
$\bar H=\tfrac12(\log 2+\log 8)=\log 4$ và $\PPL=4$.

\subsubsection{Quy tắc tổng hợp và phân biệt với DPO}
\label{sec:b6-agg}

Khi cộng qua nhiều batch hoặc nhiều câu: cộng tổng log-xác suất âm rồi chia tổng số
token được chấm; \emph{không} lấy trung bình các PPL riêng lẻ. Lý do: $\PPL$ là hàm
mũ lồi của $\bar H$, nên trung bình các PPL không bằng PPL của trung bình (bất đẳng
thức Jensen), và các batch nói chung có số token khác nhau. Loại các vị trí padding
khỏi cả tử và mẫu. Khi so sánh, dùng cùng tokenizer, dữ liệu và quy ước chấm token.

Ngược lại, $\log\pi_\theta(y\mid x)$ trong \eqref{eq:b5-logit} là log-xác suất
\emph{cả câu trả lời}. Theo quy tắc dây chuyền nó bằng
$\sum_{t=1}^{T}\log\pi_\theta(y_t\mid x,y_{<t})$ với $T$ là số token của đáp án
theo cùng quy ước kết thúc chuỗi cho hai mô hình. Không tính log-xác suất của prompt
cố định, và \emph{không} chia tổng này cho $T$ trong DPO chuẩn.

Lý do của sự khác biệt: PPL đo \emph{chất lượng trên mỗi token} nên phải chuẩn hóa
theo độ dài để so sánh được giữa các tập; còn $\log\pi_\theta(y\mid x)$ trong DPO
là xác suất của \emph{một lựa chọn} trong mô hình Bradley--Terry
\eqref{eq:b5-bt}, và chia cho $T$ sẽ thay đổi bài toán so sánh ưu tiên theo độ dài
đáp án.

\subsubsection{Đối chiếu với cài đặt A.2}

Quy tắc ở Mục~\ref{sec:b6-agg} khớp \texttt{validation\_metrics} trong
\texttt{src/char\_lm.py}: mã tích lũy \texttt{nll\_sum} ở \texttt{float64} với
\texttt{reduction='sum'}, chia cho \texttt{token\_count} rồi mới lấy
\texttt{math.exp}. Docstring ghi rõ ``không lấy trung bình PPL các batch''.

\clearpage
\section{Phần C --- Bài tập tính toán và phân tích}
\label{sec:phan-c}

Phần~C nối lý thuyết ở Phần~B với mã ở Phần~A: mỗi bài được giải bằng tay hoặc
bằng công thức tường minh, rồi đối chiếu số với cài đặt của nhóm trên CPU
\cite[tr.~19]{de}.

\textbf{Về đầu vào chính thức.} Đề ghi C.1 dùng ``ma trận Q/K/V do giảng viên cung
cấp''. Theo thông tin nhóm nhận được, giảng viên không phát bộ ma trận riêng; C.4
nêu ``cho log-xác suất'' nhưng không kèm số. Nhóm vì vậy dùng \emph{bộ minh họa có
ghi nguồn và tái lập được}, và ghi rõ đây không phải số liệu chính thức. Nếu nhận
được đầu vào riêng, chỉ cần thay số ở ô đầu vào của notebook; toàn bộ phép tính và
phần đối chiếu giữ nguyên.

%--------------------------------------------------------------------
\subsection{C.1 --- Tính tay một đầu attention có mặt nạ nhân quả}
\label{sec:c1}

\textbf{Đề bài.} Ba token, một đầu chú ý, $d_k=2$, có mặt nạ nhân quả.

\textbf{Đầu vào minh họa.} Sinh bằng \texttt{random.Random(577893521)}, lấy số
nguyên trong $[-2,2]$ theo hàng, lần lượt $Q$, $K$, $V$; không lấy mẫu lại. Mỗi
hàng là một token; chọn $d_v=2$.
\[
  Q=\begin{bmatrix}-2&0\\-1&-2\\1&-2\end{bmatrix},\quad
  K=\begin{bmatrix}-1&-1\\2&1\\0&2\end{bmatrix},\quad
  V=\begin{bmatrix}-1&0\\-1&1\\0&-1\end{bmatrix}.
\]

\subsubsection{Bước 1: điểm số, chia tỷ lệ, che tương lai}

Phần tử $(i,j)$ của $QK^\top$ là tích vô hướng của query token $i$ với key token
$j$. Ví dụ hàng 2: $(-1)(-1)+(-2)(-1)=3$; $(-1)2+(-2)1=-4$; $(-1)0+(-2)2=-4$.
Đặt $c=1/\sqrt{2}$:
\[
  QK^\top=\begin{bmatrix}2&-4&0\\3&-4&-4\\1&0&-4\end{bmatrix},
  \qquad
  S=c\,QK^\top+M=\begin{bmatrix}
      2c&-\infty&-\infty\\
      3c&-4c&-\infty\\
      c&0&-4c\end{bmatrix}.
\]
Mặt nạ cộng $M_{ij}=0$ nếu $j\le i$ và $-\infty$ nếu $j>i$. Việc che xảy ra
\emph{trước} khi chuẩn hóa từng hàng: vì $e^{-\infty}=0$, token tương lai không góp
vào cả tử \emph{và} mẫu. Đường chéo được giữ, nên mỗi hàng luôn còn ít nhất một
khóa hợp lệ.

\subsubsection{Bước 2: softmax từng hàng và tổng có trọng số}

Gọi $A=\softmax_{\mathrm{hàng}}(S)$ và $o_i=\sum_{j\le i}A_{ij}v_j$. Trừ điểm lớn
nhất của hàng trước khi lấy mũ không làm đổi softmax (tính bất biến dịch), và
tránh tràn số.

\begin{description}[leftmargin=2em,style=nextline]
\item[Token 1] Chỉ $v_1$ hợp lệ, nên $A_{1,:}=(1,0,0)$ và $o_1=v_1=(-1,0)$.
\item[Token 2] Trừ $3c$ khỏi hàng 2. Đặt $t=e^{-7c}$, khi đó
$A_{2,:}=(1,t,0)/(1+t)$ và
$o_2=\bigl[(-1,0)+t(-1,1)\bigr]/(1+t)=\bigl(-1,\ t/(1+t)\bigr)$.
Tọa độ đầu vẫn đúng bằng $-1$ vì cả hai value hợp lệ đều có tọa độ đầu bằng $-1$.
\item[Token 3] Trừ $c$ khỏi hàng 3. Đặt $u=e^{-c}$, $v=e^{-5c}$, $D=1+u+v$. Khi đó
$A_{3,:}=(1,u,v)/D$ và
$o_3=\bigl[(-1,0)+u(-1,1)+v(0,-1)\bigr]/D=\bigl(-(1+u)/D,\ (u-v)/D\bigr)$.
\end{description}

\begin{center}
\begin{tabular}{crr}
\toprule
\textbf{Token} & \textbf{Trọng số trên key 1, 2, 3} & \textbf{Đầu ra hai chiều}\\
\midrule
1 & $(1,\ 0,\ 0)$ & $(-1,\ 0)$\\
2 & $(0{,}992965,\ 0{,}007035,\ 0)$ & $(-1,\ 0{,}007035)$\\
3 & $(0{,}656939,\ 0{,}323916,\ 0{,}019145)$ & $(-0{,}980855,\ 0{,}304771)$\\
\bottomrule
\end{tabular}
\end{center}

Bảng được làm tròn \emph{sau} khi tính $O=AV$; không dùng trọng số đã làm tròn để
tính tiếp.

\subsubsection{Đối chiếu \dalam}

Notebook C.1 chạy \texttt{scaled\_dot\_product\_attention} và \texttt{causal\_mask}
của A.1 ở \texttt{float64} trên CPU, so với biểu thức đóng ở trên bằng
\texttt{torch.testing.assert\_close} với \texttt{atol=1e-12}, và kiểm thêm
\texttt{weights.triu(1).count\_nonzero() == 0} --- tức không có trọng số dương nào
ở phần tam giác trên, xác nhận không rò rỉ tương lai.

\textbf{Nhận xét đáng nêu ở vấn đáp.} Trọng số của token 2 đặt $99{,}3\%$ khối
lượng vào key 1 vì chênh lệch điểm là $7c\approx4{,}95$ --- đúng hiện tượng bão hòa
mô tả ở \eqref{eq:b1-sat} với $m=2$, $g=7c$. Đây là ví dụ cụ thể cho thấy vì sao
hệ số $1/\sqrt{d_k}$ có ý nghĩa.

%--------------------------------------------------------------------
\subsection{C.2 --- Đếm tham số và phép tính của một khối Transformer}
\label{sec:c2}

\textbf{Phạm vi.} Một lần forward của \emph{một khối} trên một chuỗi
$X\in\R^{n\times d}$, với $d=\dmodel$, $f=d_{ff}$, $h$ đầu, $h\mid d$,
$d_h=d/h$. Không tính embedding, PE, LayerNorm cuối mô hình, đầu đọc từ vựng,
backward hay optimizer. Mặt nạ nhân quả là buffer cố định, không phải tham số.

\subsubsection{Số tham số}

\texttt{Linear(a,b)} học $ab$ trọng số và $b$ bias. LayerNorm trên $d$ chiều học cả
$\gamma$ và $\beta$, mỗi vector dài $d$. Phép reshape chia đầu chỉ tổ chức lại cùng
$d$ chiều, \emph{không} nhân số tham số lên $h$ lần.

\begin{center}
\begin{tabular}{lrrr}
\toprule
\textbf{Thành phần} & \textbf{Trọng số} & \textbf{Bias} & \textbf{Tổng}\\
\midrule
Chiếu Q,K,V gộp ($d\to3d$) & $3d^2$ & $3d$ & $3d^2+3d$\\
Chiếu đầu ra ($d\to d$) & $d^2$ & $d$ & $d^2+d$\\
FFN lớp 1 ($d\to f$) & $df$ & $f$ & $df+f$\\
FFN lớp 2 ($f\to d$) & $df$ & $d$ & $df+d$\\
Hai LayerNorm & $2d$ & $2d$ & $4d$\\
\midrule
\multicolumn{3}{r}{\textbf{Tổng}} & $\mathbf{4d^2+2df+9d+f}$\\
\bottomrule
\end{tabular}
\end{center}

\begin{equation}
  N_{\mathrm{params}}=4d^2+2df+9d+f.
  \label{eq:c2-params}
\end{equation}

Với $d,f$ cố định, số tham số \emph{không} phụ thuộc $n$ hay $h$. ReLU, softmax,
phép cộng residual và mặt nạ không có tham số học.

\subsubsection{Số phép tính forward}

Quy ước MAC: tích $(a\times b)(b\times c)$ dùng $abc$ lượt nhân--tích lũy; khi đối
chiếu profiler dùng 1 MAC $=$ 2 FLOP theo quy ước GEMM.

\begin{center}
\begin{tabular}{llr}
\toprule
\textbf{Thành phần} & \textbf{Kích thước tích} & \textbf{MAC cả khối}\\
\midrule
Chiếu Q,K,V & $(n\times d)(d\times 3d)$ & $3nd^2$\\
Điểm $QK^\top$ & mỗi đầu $(n\times d_h)(d_h\times n)$ & $hn^2d_h=n^2d$\\
Tổng có trọng số $PV$ & mỗi đầu $(n\times n)(n\times d_h)$ & $hn^2d_h=n^2d$\\
Chiếu đầu ra & $(n\times d)(d\times d)$ & $nd^2$\\
FFN hai lớp & $(n\times d)(d\times f)$, $(n\times f)(f\times d)$ & $2ndf$\\
\midrule
\multicolumn{2}{r}{\textbf{Tổng}} & $\mathbf{4nd^2+2ndf+2n^2d}$\\
\bottomrule
\end{tabular}
\end{center}

\begin{equation}
  \mathrm{MAC}_{\mathrm{matmul}}=4nd^2+2ndf+2n^2d.
  \label{eq:c2-mac}
\end{equation}

Lưu ý phép triệt tiêu $h\,d_h=d$: chi phí attention \emph{không} phụ thuộc số đầu
khi $d$ giữ nguyên. Đây là điểm bị hỏi thường xuyên.

\textbf{Ngoài các tích ma trận:} cộng bias $n(5d+f)$; hai residual $2nd$; ReLU
$\Theta(nf)$; hai LayerNorm $\Theta(nd)$; chia tỷ lệ điểm $hn^2$ phép nhân và
softmax xử lý $hn^2$ phần tử, chi phí $\Theta(hn^2)$.

\subsubsection{Hạng tử bậc hai theo $n$}

Trong \eqref{eq:c2-mac}, hai hạng $n^2d$ đến từ $QK^\top$ và $PV$; phần chiếu và
FFN tăng \emph{tuyến tính} theo $n$. Vì $h\le d$, độ phức tạp toàn khối là
$\Theta(nd^2+ndf+n^2d)$, và score cùng xác suất attention có kích thước
$h\times n\times n$ --- tức bộ nhớ cũng bậc hai. Với batch $B$, số phép tính nhân
$B$ còn số tham số giữ nguyên.

\textbf{Vì sao mặt nạ nhân quả không chia đôi chi phí?} Cài đặt A.1 tính đủ ma trận
$n\times n$ rồi mới điền $-\infty$ phía trên đường chéo, và $PV$ vẫn là phép nhân
dense đủ kích thước dù các trọng số bị che bằng $0$. Chỉ thuật toán \emph{thực sự
bỏ qua} ô bị che mới khai thác được $n(n+1)/2$ cặp hợp lệ; mặt nạ trong A.1 không
làm việc đó. Đây là khác biệt giữa đúng về ngữ nghĩa và tối ưu về chi phí.

\subsubsection{Thay số theo cấu hình A.2 \dalam}

Với $d=64$, $f=128$, $h=4$, $n=128$:
\begin{align*}
  N_{\mathrm{params}}&=4(64)^2+2(64)(128)+9(64)+128\\
   &=16{.}384+16{.}384+576+128=\mathbf{33{.}472},\\
  \mathrm{MAC}&=4(128)(64)^2+2(128)(64)(128)+2(128)^2(64)\\
   &=2{.}097{.}152+2{.}097{.}152+2{.}097{.}152=\mathbf{6{.}291{.}456}.
\end{align*}
Tương đương $12{.}582{.}912$ FLOP theo quy ước trên. Ở cấu hình này hai tích
attention chiếm \emph{đúng} $1/3$ tổng MAC; tỷ lệ đó tăng khi chuỗi dài hơn.

Notebook C.2 đối chiếu hai chiều: \texttt{numel()} của từng submodule khớp
\eqref{eq:c2-params} cho cả $h=4$ và $h=8$ (xác nhận tính độc lập với $h$), và FLOP
mà PyTorch Profiler ước lượng cho \texttt{aten::addmm} (chiếu, FFN) cùng
\texttt{aten::bmm} (hai tích attention) khớp đúng $2\times$ \eqref{eq:c2-mac} với
$n\in\{64,128,256\}$. Khi $n$ gấp đôi, chiếu và FFN gấp đôi còn hai tích attention
gấp bốn --- xác nhận hạng tử $O(n^2)$ bằng số.

Không dùng số FLOP của một khối forward làm FLOP huấn luyện toàn mô hình.

%--------------------------------------------------------------------
\subsection{C.3 --- Tham số LoRA ở Q/V và bộ nhớ QLoRA 4 bit}
\label{sec:c3}

\textbf{Mô hình nền dùng để tính.} Qwen2.5-1.5B-Instruct, khóa theo revision
\texttt{989aa7980e4cf806f80c7fef2b1adb7bc71aa306} \cite{qwen}. Đây là mô hình dùng
cho \emph{phép tính} C.3; nó không đồng nghĩa với quyết định mô hình của D.3.

Kiến trúc: $L=28$ khối, chiều ẩn $d=1536$, 12 đầu Q nhưng chỉ 2 đầu K/V
(grouped-query attention), mỗi đầu 128 chiều. Do đó trọng số Q có dạng
$1536\times1536$ còn V là $256\times1536$ (đầu ra $\times$ đầu vào). Việc lặp K/V
khi tính attention \emph{không} tạo thêm trọng số.

\subsubsection{Đếm tham số học}

Theo \cite[§4.1]{lora}, với $W$ kích thước $o\times i$, hai ma trận học có kích
thước $r\times i$ và $o\times r$, tức $r(i+o)$ tham số --- đúng
\eqref{eq:b4-count}. Đóng băng toàn bộ mô hình nền kể cả bias; chỉ gắn adapter vào
\texttt{q\_proj} và \texttt{v\_proj} ở mỗi khối, không học K, O, embedding hay đầu
đọc. Hệ số $\alpha/r$ không làm đổi số tham số.
\begin{align*}
  P_Q&=28(1536+1536)r=86{.}016\,r,\\
  P_V&=28(1536+256)r=50{.}176\,r,\\
  P&=136{.}192\,r.
\end{align*}

\begin{center}
\begin{tabular}{rrrr}
\toprule
$r$ & \textbf{Q} & \textbf{V} & \textbf{Tổng}\\
\midrule
4 & 344.064 & 200.704 & \textbf{544.768}\\
8 & 688.128 & 401.408 & \textbf{1.089.536}\\
16 & 1.376.256 & 802.816 & \textbf{2.179.072}\\
\bottomrule
\end{tabular}
\end{center}

Với $r=8$, tỷ lệ tham số học so với $1{,}54$ tỷ tham số của mô hình nền là khoảng
$0{,}071\%$.

Notebook C.3 kiểm số này bằng chính lớp \texttt{LoRALinear} của A.3 trên
\emph{meta tensor} --- tức kiểm kích thước mà không cấp vùng nhớ, không lượng tử
hóa và không tải Qwen.

\subsubsection{Ước lượng bộ nhớ khi kết hợp QLoRA 4 bit}

\textbf{Kịch bản:} SFT, không offload, không có reference policy của DPO. Phân biệt
rõ \emph{trọng số lưu} NF4 với \emph{kiểu tính toán} 16 bit.

Bảy ma trận lớn mỗi khối (Q/K/V/O và gate/up/down) lưu NF4 với double
quantization. Gọi $k=256$ là chiều ra K/V và $f=8960$ là chiều trung gian. Số phần
tử lượng tử hóa là
\[
  N_4=L\,(2d^2+2dk+3df).
\]
Theo \cite[§3, tr.~5]{qlora}, mỗi phần tử dùng trung bình
$4+\tfrac{8}{64}+\tfrac{32}{64\cdot256}$ bit, kể cả scale hai cấp. Bỏ qua codebook,
offset và padding theo tensor.

Giả sử embedding và đầu đọc dùng chung trọng số
(\texttt{tie\_word\_embeddings=true}), RMSNorm và bias lưu 16 bit:
\[
  N_{16}=V_{\mathrm{ocab}}\,d+(2L+1)d+L(d+2k),
\]
gồm một embedding, hai norm mỗi khối, một norm cuối, và bias Q/K/V. Không đếm
\texttt{lm\_head} lần thứ hai. Đây là \emph{giả định lưu trữ}, không phải mặc định
của mọi thư viện: nếu chuyển phần này sang FP32 thì cộng thêm $2N_{16}$ byte.

Adapter lưu FP32, gradient FP32, hai moment AdamW FP32 thường trú, tức
\textbf{16 byte cho mỗi tham số học}; không tính gradient hay optimizer cho
backbone đóng băng. Sau khi trạng thái optimizer đã cấp phát và gradient còn sống:
\begin{equation}
  M_{\mathrm{tĩnh}}=\frac{N_4\,b_4}{8}+2N_{16}+16P \quad\text{byte}.
  \label{eq:c3-static}
\end{equation}

Notebook C.3 có một đối chứng độc lập đáng chú ý: $N_4+N_{16}=1{.}543{.}714{.}304$,
khớp đúng tổng số phần tử BF16 của checkpoint trên Hub trước khi lượng tử hóa. Đây
là phép kiểm bắt được lỗi đếm thiếu hoặc đếm trùng một họ ma trận.

\subsubsection{Từ bộ nhớ tĩnh tới kịch bản có activation}

Chọn microbatch $B=1$, chiều dài $T=512$, checkpoint từng khối, không KV cache.
Ước lượng các tensor chính ở pha tính loss: một đầu vào FP16 cho mỗi khối
($2LBTd$ byte), cùng logits FP32 và một tensor FP32 cùng cỡ cho loss không fused
($2\times4BTV_{\mathrm{ocab}}$ byte).

\textbf{Kết luận định tính.} Tăng rank từ 4 lên 16 làm số tham số học tăng 4 lần
nhưng chỉ thêm khoảng 25 MiB trạng thái adapter theo \eqref{eq:c3-static}, không
nhân bốn bộ nhớ cả mô hình. Phần chi phối ngân sách là backbone lượng tử hóa và các
tensor từ vựng, không phải adapter.

\textbf{Giới hạn \chualam.} Đây là \emph{ước lượng theo các tensor được nêu}, không
phải phép đo. Chưa cộng: tái tính activation và backward của block, buffer giải
lượng tử, workspace GEMM/attention, CUDA context, tensor tạm của optimizer, và phân
mảnh allocator. Khi $T$ tăng, hai khoản activation tăng tuyến tính; nếu attention
lưu điểm dense thì riêng một ma trận điểm FP16 cỡ $B\times h\times T\times T$ tăng
\emph{bậc hai}. Vì vậy \textbf{không} được suy từ bảng rằng một GPU cụ thể chạy vừa
SFT hay DPO. Muốn xác nhận VRAM cho D.3 phải đo cấu hình thật qua ít nhất một bước
forward, backward và optimizer sau khi cấp phát trạng thái, ghi cả peak allocated
và peak reserved.

%--------------------------------------------------------------------
\subsection{C.4 --- Mất mát DPO, gradient và ảnh hưởng của $\beta$}
\label{sec:c4}

\subsubsection{Lời giải cho một cặp bất kỳ}

Cùng một prompt $x$, gọi $w$ là đáp án được ưa thích, $l$ là đáp án bị từ chối.
Đặt $p_i=\log\pi_\theta(y_i\mid x)$ và $r_i=\log\pref(y_i\mid x)$ với
$i\in\{w,l\}$. Đây là \emph{tổng} log-xác suất tự nhiên của token đáp án, không gồm
prompt hay padding, và không phải trung bình theo độ dài như trong perplexity
(Mục~\ref{sec:b6-agg}). Các log-xác suất hữu hạn; $\beta>0$; reference đóng băng.

Theo \eqref{eq:b5-logit} và \cite[§4, Eq. (7)]{dpo}:
\begin{equation}
  m=(p_w-p_l)-(r_w-r_l),\qquad z=\beta m,\qquad
  L=-\log\sigma(z)=\log(1+e^{-z}).
  \label{eq:c4-loss}
\end{equation}
Điểm cần nhấn: $m$ so \emph{tỷ số} xác suất của hai đáp án giữa policy và
reference, không chỉ so $p_w$ với $p_l$. Và $\sigma(z)$ là xác suất \emph{ưu tiên},
khác xác suất policy sinh ra riêng đáp án $w$.

Vì $\mathrm dL/\mathrm dz=-\sigma(-z)$ và
$\partial z/\partial p_w=\beta=-\partial z/\partial p_l$, đặt $g=\beta\sigma(-z)$:
\begin{equation}
  \frac{\partial L}{\partial p_w}=-g,\qquad
  \frac{\partial L}{\partial p_l}=+g,\qquad
  \frac{\partial L}{\partial r_w}=\frac{\partial L}{\partial r_l}=0 .
  \label{eq:c4-grad}
\end{equation}
Hai gradient đối nhau: mất mát đẩy log-xác suất của $w$ lên và của $l$ xuống với
cùng độ lớn. Gradient theo reference bằng $0$ vì reference bị
\texttt{detach} --- phép kiểm trong notebook xác nhận điều này bằng
\texttt{assert grw is None and grl is None}.

\subsubsection{Thế số --- ví dụ minh họa}

Chọn $(p_w,p_l,r_w,r_l)=(\log 0{,}2,\ \log 0{,}1,\ \log 0{,}4,\ \log 0{,}1)$ và
$\beta=1$. Hai đáp án không phải toàn bộ không gian đầu ra, nên hai xác suất của
mỗi mô hình không cần cộng thành $1$.

Policy ưu tiên $w$ so với $l$ theo tỷ số $2:1$, nhưng reference ưu tiên mạnh hơn,
$4:1$. Vậy
\[
  m=\log 2-\log 4=-\log 2,\quad z=-\log 2,\quad \sigma(z)=\tfrac13,
\]
\[
  \boxed{L=\log 3\approx 1{,}098612},\qquad
  g=\tfrac23,\qquad
  \left(\frac{\partial L}{\partial p_w},\frac{\partial L}{\partial p_l}\right)
   =\left(-\tfrac23,\ +\tfrac23\right).
\]

\textbf{Quan sát quan trọng:} mất mát $\log 3$ \emph{lớn hơn} $\log 2$ dù
$p_w>p_l$. Lý do: policy đã ưu tiên $w$, nhưng chưa ưu tiên mạnh \emph{bằng}
reference, nên $m<0$. Đây là điểm dễ bị hỏi vặn nhất ở C.4, và nó chính là hệ quả
của việc \eqref{eq:b5-logit} so hai log-tỷ số chứ không so hai log-xác suất.

\subsubsection{Ảnh hưởng của $\beta$ tại một cặp cố định}

Ở một cặp cố định, $\partial L/\partial\beta=-m\,\sigma(-\beta m)$. Khi
$\beta\to0^+$, mất mát tiến tới $\log 2$ và $g\sim\beta/2$: tín hiệu gradient yếu
cho \emph{mọi} cặp.

\begin{center}
\begin{tabular}{p{0.32\textwidth}p{0.28\textwidth}p{0.26\textwidth}}
\toprule
\textbf{Policy so với reference} & \textbf{Khi tăng $\beta$} & \textbf{Khi $\beta$ rất lớn}\\
\midrule
$m<0$: chênh lệch chưa đủ & Mất mát tăng; $|g|$ tăng & $L\sim\beta|m|$, $g\sim\beta$\\
$m=0$: hai tỷ số bằng nhau & $L=\log2$ không đổi, nhưng gradient \emph{không} bằng $0$ & $g=\beta/2$\\
$m>0$: đã vượt reference & Mất mát giảm; $g$ không tăng đơn điệu & $L\to0$, $g\to0$ do bão hòa\\
\bottomrule
\end{tabular}
\end{center}

Hàng giữa đáng chú ý: khi hai tỷ số bằng nhau, mất mát bằng $\log 2$ \emph{bất kể}
$\beta$, nhưng gradient bằng $\beta/2\ne0$ --- mô hình vẫn bị đẩy đi. Đây là khác
biệt giữa ``mất mát không đổi'' và ``không học gì''.

\textbf{Phân biệt hai cách đọc $\beta$.} Bảng trên chỉ đổi $\beta$ tại policy cố
định. Trong mục tiêu gốc $\E[R]-\beta\KL(\pi\|\pref)$ ở \eqref{eq:b5-obj}, giữ cố
định reward và reference thì $\beta$ lớn đặt trọng số cao hơn lên phạt KL. Với cùng
prompt và hai đáp án có xác suất reference dương, \eqref{eq:b5-gibbs} cho tỷ số tối
ưu
\[
  \frac{\pi^*(y_w\mid x)}{\pi^*(y_l\mid x)}
   =\frac{\pref(y_w\mid x)}{\pref(y_l\mid x)}
    \exp\!\left(\frac{R_w-R_l}{\beta}\right).
\]
Nghiệm thay đổi theo $\beta$ \emph{nếu} reward không hằng trên hỗ trợ của reference.
Nếu $R(x,y)=c(x)$ trên hỗ trợ đó thì hệ số $e^{c(x)/\beta}$ là chung cho mọi đáp án,
bị triệt tiêu khi chuẩn hóa, và $\pi^*=\pref$ với mọi $\beta>0$. Hai cách đọc này
không được lẫn.

\subsubsection{Đối chiếu \dalam}

Notebook C.4 kiểm trên CPU ở \texttt{float64}, dùng hai dạng ổn định số cho
$\mathrm{softplus}(-z)$ và $\sigma(-z)$, so với \texttt{src/dpo.py} và với
\texttt{torch.autograd.grad}: mất mát và cả hai gradient policy khớp ở
\texttt{rel\_tol=1e-12}; gradient theo reference xác nhận là \texttt{None}. Phép
kiểm chạy qua \textbf{ba dấu của margin} ($m=-\log2,\ 0,\ +\log2$, bằng cách đổi
xác suất policy cho $w$ thành $0{,}2$, $0{,}4$, $0{,}8$) và \textbf{bốn giá trị}
$\beta\in\{0{,}1;\ 1;\ 2;\ 8\}$. Mỗi hàng là một cặp riêng; không lấy trung bình
giữa các ví dụ.

\textbf{Giới hạn.} Bộ số là minh họa, không phải số giảng viên cấp, không phải đầu
ra của một checkpoint, và không phải nhãn người chấm ở D.4. Không dùng bảng này để
chọn $\beta$ cho D.4 hay kết luận DPO cải thiện chất lượng.

%--------------------------------------------------------------------
\subsection{Tổng hợp trạng thái Phần C}

\begin{center}
\begin{tabular}{llp{0.42\textwidth}}
\toprule
\textbf{Mã} & \textbf{Trạng thái} & \textbf{Kết quả chính}\\
\midrule
C.1 & \dalam & Ba token giải từng hàng, khớp A.1 ở \texttt{atol=1e-12}; đầu vào minh họa seed \texttt{577893521}\\
C.2 & \dalam & 33.472 tham số, 6.291.456 MAC tại $n=128$; hạng tử $O(n^2)$ xác nhận bằng profiler\\
C.3 & \dalam & 544.768 / 1.089.536 / 2.179.072 tham số học; ước lượng QLoRA, chưa đo VRAM\\
C.4 & \dalam & $L=\log3$, gradient $(-2/3,+2/3)$; khớp A.4 và autograd trên 3 dấu margin $\times$ 4 giá trị $\beta$\\
\bottomrule
\end{tabular}
\end{center}

\clearpage
\section{Phần D --- Ứng dụng: trợ lý thủ tục hành chính}
\label{sec:phan-d}

Phần~D là bài toán Việt Nam, trọng số 30 điểm, gồm năm hạng mục D.1--D.5
\cite[tr.~19]{de}. Mục này trình bày \emph{trạng thái thực tế}: những gì đã làm có
bằng chứng trong kho mã, và những ngưỡng của đề chưa đạt được. Nhóm không trình
bày kế hoạch như kết quả.

\textbf{Tóm tắt trạng thái.} D.1 \motphan; D.2 \motphan\ (400 QA ứng viên, 24 đã
duyệt); D.3 \motphan\ (pilot 24 QA đã chạy trên Kaggle, chưa chạy tập chính);
D.4 \chualam; D.5 \chualam.

%--------------------------------------------------------------------
\subsection{D.1 --- Thu thập mô tả thủ tục công khai}
\label{sec:d1}

\textbf{Yêu cầu.} Thu thập trình tự, hồ sơ, thời hạn, lệ phí và căn cứ pháp lý từ
Cổng Dịch vụ công quốc gia và cổng địa phương; lưu truy xuất và phiên bản.

\subsubsection{Năm họ thủ tục đã rà \motphan}

\begin{center}
\small
\begin{tabular}{p{0.40\textwidth}lp{0.28\textwidth}}
\toprule
\textbf{Họ thủ tục} & \textbf{Mã} & \textbf{Trạng thái căn cứ}\\
\midrule
Cấp bản sao trích lục hộ tịch / bản sao giấy khai sinh & \texttt{2.000635} &
Đã xác minh QĐ 1818 và QĐ 2163; \emph{thời hạn chưa thống nhất}\\
Chứng thực bản sao từ bản chính & \texttt{2.000815} &
Đã xác minh Điều 20 NĐ 23/2015 hợp nhất 753/VBHN-BTP\\
Chứng thực chữ ký & \texttt{2.000884} &
Đã xác minh phụ lục QĐ 858; chưa chốt quy trình tại Hà Nội\\
Thay đổi, cải chính, bổ sung hộ tịch, xác định lại dân tộc & \texttt{1.004859} &
Đã xác minh QĐ 163 và biểu lệ phí NQ 98/2026/NQ-HĐND\\
Xác nhận thông tin hộ tịch & \texttt{2.002516} &
Đã xác minh QĐ 1852; chuyển dẫn chứng từ video sang văn bản\\
\bottomrule
\end{tabular}
\end{center}

Mỗi bản ghi lưu URL, thời điểm truy xuất và SHA-256 của bản tải. Nhóm cũng kiểm
chữ ký CMS về toàn vẹn nội dung trên bốn PDF Hà Nội --- đây là kiểm toàn vẹn tệp,
\emph{không} phải thẩm định chuỗi tin cậy.

\subsubsection{Hai khác biệt phải xử lý trước khi huấn luyện}

Quá trình rà phát hiện hai mâu thuẫn mà nếu bỏ qua sẽ sinh dữ liệu sai:

\begin{enumerate}[leftmargin=2em]
\item \textbf{Thời hạn cấp trích lục không thống nhất giữa các văn bản.} Phụ lục
QĐ 1818/QĐ-TTPVHCC (hiệu lực 01/01/2026) ghi \emph{dự kiến 3 ngày làm việc}, tối đa
10 ngày nếu phải xác minh. QĐ 2163/QĐ-BTP (30/09/2026) ghi \emph{trong ngày tiếp
nhận, hoặc ngày làm việc kế tiếp nếu sau 15 giờ}. Nhóm \emph{chưa} tìm được văn bản
triển khai 2163 tại Hà Nội, cũng chưa tìm được văn bản xác nhận 1818 đã bị thay
thế. Xử lý: không dùng ``3/10 ngày'' làm đáp án chung hiện hành; chỉ soạn câu hỏi
gắn với \emph{từng văn bản và phạm vi của nó}.
\item \textbf{Lệ phí không phải con số hiển thị trên cổng.} Thông tư
226/2016/TT-BTC là căn cứ phí \emph{chứng thực}, không phải biểu lệ phí hộ tịch Hà
Nội; biểu hiện hành của Hà Nội nằm ở Nghị quyết 98/2026/NQ-HĐND. Trộn hai khái niệm
này tạo ra đáp án sai một cách khó phát hiện.
\end{enumerate}

Hai phát hiện này minh họa vì sao hạng mục D.1 không thể rút gọn thành ``cào trang
thủ tục'': nội dung hiển thị trên cổng không đồng nhất với văn bản có hiệu lực.

\subsubsection{Căn cứ quyền sử dụng}

Nhóm phân loại nội dung trước khi đưa vào dữ liệu:

\begin{itemize}[leftmargin=2em]
\item \textbf{Chính văn quy phạm và văn bản hành chính do cơ quan ban hành:} theo
khoản 2 Điều 15 Luật Sở hữu trí tuệ (VBHN 160/2026) và khoản 2 Điều 8 VBHN
6154/2026, nhóm văn bản này \emph{không} thuộc phạm vi bảo hộ quyền tác giả. Có
căn cứ để dùng phần chính văn công khai đã xác minh làm nguồn soạn QA, kèm số hiệu,
điều/mục, phiên bản và URL.
\item \textbf{Trang chi tiết của cổng, bài tin, bài hướng dẫn, video, bản lưu
pilot:} dùng để \emph{tìm} văn bản và đối chiếu cách hiển thị. Điều khoản Cổng quốc
gia yêu cầu dùng lại đúng pháp luật; footer cổng Hà Nội ghi
\emph{All Rights Reserved} và yêu cầu ghi nguồn. Nhóm \textbf{chưa} chốt căn cứ để
lấy nguyên văn toàn trang hoặc phân phối lại chúng, nên khi cần cùng một quy định
thì trích từ văn bản chính thức, không chuyển snapshot web sang tập train.
\end{itemize}

Về truy cập: \texttt{robots.txt} của cổng Hà Nội cho \texttt{Allow: /} với
\texttt{Crawl-delay: 5}, nên nhóm truy xuất tuần tự cách ít nhất 5 giây và không
crawl toàn cổng. Với các host trả HTTP 404 cho \texttt{robots.txt}, nhóm ghi rõ
rằng \emph{không có tệp robots không đồng nghĩa có giấy phép nội dung}.

\textbf{Còn thiếu.} Chưa phát hành bộ nguồn huấn luyện; còn phải kiểm các kênh và
phiên bản triển khai thực tế tại Hà Nội.

%--------------------------------------------------------------------
\subsection{D.2 --- Bộ dữ liệu chỉ dẫn có trích dẫn}
\label{sec:d2}

\textbf{Yêu cầu của đề.} Tạo $\ge 2.000$ QA có trích dẫn bằng phương pháp bán tự
động; kiểm thủ công $\ge20\%$; tập test là câu hỏi của \emph{người dùng thực} thu
qua khảo sát, tách hoàn toàn khỏi train.

\subsubsection{Đã làm \motphan}

\begin{itemize}[leftmargin=2em]
\item \textbf{Quy trình sinh có cổng kiểm.} Nhóm viết recipe soạn QA từ chính văn,
kiểm SHA-256 và từng đoạn trích, CLI chọn mẫu phân tầng, và một cổng xuất SFT từ
chối phát hành khi dữ liệu chưa đạt điều kiện.
\item \textbf{Công cụ chặn rò rỉ test.} Kiểm tách \texttt{USER\_TEST}, chặn đường
dẫn thiếu, trùng họ thủ tục và trùng câu.
\item \textbf{Lô pilot chính văn v2 đã duyệt và phát hành nội bộ:} 24 QA, nhóm
kiểm thủ công \textbf{24/24}, có biên bản chung của lô. Cổng D.2 trả
\texttt{pilot\_sft\_ready}; đầu vào D.3 kiểm \texttt{release-valid} với 18 train
và 6 validation. Đây là lô duy nhất đã phát hành.
\item \textbf{Lô pilot v1 trước đó:} 20 QA, 21 đoạn trích, 9 nguồn; đã rà riêng.
Validator ngoại tuyến xác nhận 20/20 qua kiểm cấu trúc và 21/21 đoạn trích khớp
bản lưu --- nhưng kiểm tự động không được tính vào quota kiểm thủ công.
\item \textbf{Tập chính đang mở rộng theo từng lô:} b01 thêm 57 QA train, b02
thêm 54, b03 thêm 265. Tích lũy \textbf{394 QA train và 6 validation ứng viên},
\emph{chưa phát hành}. Các lô mới gồm cả câu cần từ chối hoặc nói rõ giới hạn
nguồn.
\end{itemize}

\subsubsection{Khoảng cách tới ngưỡng của đề}

\begin{center}
\begin{tabular}{lrrl}
\toprule
\textbf{Ngưỡng} & \textbf{Yêu cầu} & \textbf{Hiện có} & \textbf{Trạng thái}\\
\midrule
Số QA có trích dẫn & $\ge2.000$ & 400 ứng viên & \chualam\\
Kiểm thủ công & $\ge20\%$ ($\ge400$) & 24 đã duyệt & \chualam\\
Tập test từ khảo sát người thật & có & chưa có & \chualam\\
Tách train/validation/test & có & có cho lô pilot & \motphan\\
Datasheet & có & đã soạn khung & \motphan\\
\bottomrule
\end{tabular}
\end{center}

\textbf{Vì sao tập test không thể sinh bằng máy.} Đề yêu cầu test là câu hỏi của
người dùng thực. Mục đích là đo khả năng tổng quát hóa sang \emph{cách người dân
thật diễn đạt} --- thường thiếu thuật ngữ, thiếu bối cảnh, hoặc mô tả tình huống cá
nhân thay vì hỏi tên thủ tục. Dùng AI sinh câu hỏi thay khảo sát sẽ phá chính phép
đo đó, vì câu hỏi sinh ra mang phân phối của mô hình, không phải của người dùng.
Quy tắc nội bộ của nhóm cấm việc này.

Thu dữ liệu người cần văn bản đồng ý và ẩn danh trước khi lưu
(Mục~\ref{sec:dao-duc}), nên hạng mục này có một \emph{sàn thời gian} không nén
được bằng công sức.

%--------------------------------------------------------------------
\subsection{D.3 --- Tinh chỉnh có giám sát bằng QLoRA}
\label{sec:d3}

\textbf{Yêu cầu.} Chọn LM mã nguồn mở hỗ trợ tiếng Việt, 0,5--3B tham số, giấy phép
cho nghiên cứu; SFT bằng LoRA hoặc QLoRA; tối ưu $r$, $\alpha$, learning rate,
epochs \emph{chỉ trên validation}.

\subsubsection{Đã chuẩn bị \motphan}

\begin{itemize}[leftmargin=2em]
\item \textbf{Mô hình:} Qwen2.5-1.5B-Instruct, Apache-2.0, revision được khóa
trong \texttt{scripts/}. Nằm trong khoảng 0,5--3B và hỗ trợ tiếng Việt. Giấy phép
và revision đã đối chiếu với model card \cite{qwen}.
\item \textbf{Mã:} bộ nạp dữ liệu có che prompt, mã QLoRA, notebook Kaggle, script
đóng gói tối thiểu, và kế hoạch so cấu hình $r$, $\alpha$, LR, epoch trên
validation.
\item \textbf{Cổng an toàn đã kiểm:} ZIP chỉ chứa mã thì không train; sai hash thì
dừng; CLI từ chối manifest \texttt{candidate\_only} của lô 24 QA. Nghĩa là quy
trình \emph{không thể} vô tình huấn luyện trên dữ liệu chưa được duyệt.
\item \textbf{Kiểm phần mềm trên CPU:} 7 test bộ nạp và mặt nạ đều qua. Với một câu
minh họa dùng tokenizer Qwen đúng revision: 71 token prompt bị che, giải mã 19
token đích khớp nguyên đáp án có dấu và trích dẫn, đúng một EOS. Đã chạy forward
trên một Qwen thu nhỏ với trọng số khởi tạo và adapter Q/V.
\item \textbf{Kiểm resume:} với mô hình thu nhỏ, checkpoint khôi phục đúng trọng số
adapter, trạng thái AdamW, RNG CPU và thứ tự trộn.
\end{itemize}

\subsubsection{Kết quả pilot đã chạy trên Kaggle \dalam}

Nhóm đã chạy một lượt SFT pilot thật trên lô 24 QA đã phát hành. Số liệu dưới đây
đối chiếu từ \texttt{metrics.jsonl} và \texttt{summary.json} do nhóm gửi.

\begin{center}
\begin{tabular}{lrr}
\toprule
\textbf{Epoch} & \textbf{Validation NLL} & \textbf{Validation PPL}\\
\midrule
1 & 1,509630 & 4,525054\\
2 & \textbf{1,242314} & \textbf{3,463617}\\
\bottomrule
\end{tabular}
\end{center}

Seed 41 hoàn thành $2/2$ epoch, không dừng sớm; mỗi epoch giám sát 1.235 token
train và 388 token validation. Epoch 2 được chọn theo validation NLL. PPL khớp
$e^{\mathrm{NLL}}$ theo \eqref{eq:b6-ppl}.

\textbf{Cấu hình đã chạy:} QLoRA NF4 double quantization, FP16, adapter ở Q/V,
$r=8$, $\alpha=16$, dropout $0$, learning rate $10^{-4}$, batch $1$, độ dài tối đa
$512$. Môi trường: Tesla T4, Python 3.13.16, torch 2.11.0+cu130, CUDA 13.0.

\textbf{Bộ nhớ đã đo:} peak allocated $2{,}402$ GiB, peak reserved $2{,}725$ GiB
theo PyTorch. Đây là số của tiến trình, không phải tổng VRAM toàn thiết bị. Con số
này nay \emph{thay thế} phần ước lượng ở Mục~\ref{sec:c3} cho đúng cấu hình trên.

\textbf{Tái lập đã kiểm:} bốn SHA-256 (manifest release-01, hai tệp Python trong
ZIP mã, chat template tokenizer) đều khớp. Resume từ \texttt{epoch-0001} tái lập
đúng epoch 2: cùng 1.235 token train (train NLL $1{,}750317$), cùng 388 token
validation, NLL $1{,}242314$ và PPL $3{,}463617$; chỉ hai trường peak VRAM lệch.
Resume \emph{không} được tính là một seed độc lập.

%--------------------------------------------------------------------
\subsubsection{Hỏi thử adapter pilot: hai lỗi quan sát được}

Nhóm chạy một notebook hỏi thử riêng, suy luận từ \texttt{best\_adapter/} và
\emph{không} huấn luyện lại. Mô hình sinh được hai đáp án, nhưng \textbf{cả hai
đều không đạt} mục tiêu đặt ra ở Mục~\ref{sec:gioi-thieu}:

\begin{enumerate}[leftmargin=2em]
\item \textbf{Sai thời hạn và tự nhận là mới nhất.} Với câu hỏi về xác nhận
thông tin hộ tịch, mô hình trả ``15/30 ngày làm việc'' và tự nhận đây là quy định
mới nhất. Phụ lục 1852/QĐ-TTPVHCC, B.4 tr.~3 ghi \emph{dự kiến 03 ngày}, và không
quá 10 ngày làm việc nếu cần kiểm tra, xác minh.
\item \textbf{Không từ chối câu ngoài phạm vi.} Với câu hỏi về thị thực du lịch
Nhật Bản --- hoàn toàn ngoài phạm vi hộ tịch và chứng thực tại Hà Nội --- mô hình
\emph{hướng dẫn} thay vì nêu rõ là ngoài phạm vi.
\end{enumerate}

Thêm nữa, cả hai đáp án chỉ ghi \texttt{[1]} mà không kèm nguồn truy xuất được.

\textbf{Vì sao hai lỗi này đáng ghi vào báo cáo.} Lỗi thứ hai chính là hiện tượng
mà Mục~\ref{sec:d5} đặt giả thuyết: SFT dạy mô hình \emph{luôn} trả lời, nên
``từ chối đúng cách'' là thứ SFT khó dạy. Quan sát này \emph{ủng hộ} giả thuyết
nhưng không chứng minh nó, và nhất định không phải điểm D.5.

\textbf{Đây là quan sát định tính, không phải phép đo.} Hai câu, một adapter huấn
luyện trên 24 QA, không có người đánh giá mù. Nhóm cũng \emph{không} chuyển hai QA
validation cùng dẫn đoạn B.4 sang train để làm đẹp kết quả demo --- làm vậy là rò
rỉ validation.

\subsubsection{Hạng mục D.3 còn lại \chualam}

\begin{itemize}[leftmargin=2em]
\item \textbf{Chưa chọn cấu hình.} Đề yêu cầu tối ưu $r$, $\alpha$, learning rate
và epochs trên validation. Hiện chỉ có \emph{một} cấu hình đã chạy ($r=8$,
$\alpha=16$, LR $10^{-4}$), nên chưa có việc chọn nào diễn ra.
\item \textbf{Chưa chạy trên tập chính.} Lượt pilot dùng 24 QA; 394 QA train của
tập chính chưa phát hành nên chưa huấn luyện.
\item \textbf{Chưa đủ 5 lần chạy độc lập.} Chỉ có seed 41. Yêu cầu tái lập ở
Mục~\ref{sec:tai-lap} cần ít nhất 5 seed; resume không tính là seed thứ hai.
\end{itemize}

\subsection{D.4 --- Dữ liệu ưu tiên và huấn luyện DPO}
\label{sec:d4}

\textbf{Yêu cầu.} Thu tối thiểu \textbf{500 cặp} so sánh: hai đáp án SFT cho cùng
một câu hỏi; \textbf{ít nhất hai người đánh giá độc lập} theo đúng/đầy đủ/trích
dẫn; báo cáo mức đồng thuận; huấn luyện DPO.

\textbf{Trạng thái: \chualam.}

\subsubsection{Quy trình đã thiết kế}

Nhóm đã viết hướng dẫn chấm và rubric, và chốt quy trình:

\begin{enumerate}[leftmargin=2em]
\item Với mỗi câu hỏi \emph{ngoài tập test}, sinh hai đáp án từ checkpoint SFT
\emph{thật}. Khóa checkpoint, template và cấu hình sinh trước khi chấm.
\item Hà Mạnh Thế và Trần Duy Hưng chấm \emph{độc lập} theo ba tiêu chí: đúng nội
dung, đủ ý, trích dẫn đúng. \textbf{Không chia đôi 500 cặp cho hai người} --- mỗi
người chấm toàn bộ.
\item Không ai được xem nhãn của người kia trước khi khóa.
\item Khóa nhãn, \emph{đo mức đồng thuận}, rồi Cao Văn Hà mới phân xử. Giữ nguyên
nhãn gốc và lý do loại cặp; không ghi đè nhãn raw.
\item Dùng các cặp ưu tiên để huấn luyện DPO, giữ mô hình SFT tham chiếu cố định,
và chọn $\beta$ \emph{trên validation DPO}.
\end{enumerate}

Thiết kế này trực tiếp phản ánh Mục~\ref{sec:b5}: $\pref$ phải là một mô hình xác
định và cố định, vì \eqref{eq:b5-gibbs} định nghĩa nghiệm tối ưu \emph{tương đối}
với nó; hai đáp án phải thuộc cùng một prompt để hạng $\beta\log Z(x)$ triệt tiêu
ở \eqref{eq:b5-invert}.

\subsubsection{Hai lý do hạng mục này bị chặn}

\begin{enumerate}[leftmargin=2em]
\item \textbf{Chặn kỹ thuật --- đã gỡ một phần.} Trước đây không có checkpoint SFT
nào, nên không sinh được hai đáp án để so sánh. Nay đã có adapter pilot
(\texttt{best\_adapter/}, seed 41, epoch 2), nên \emph{về mặt kỹ thuật} nhóm đã
sinh được cặp đáp án. Giới hạn còn lại là chất lượng: adapter này huấn luyện trên
24 QA và hỏi thử cho thấy nó sai thời hạn và không từ chối câu ngoài phạm vi, nên
cặp sinh từ nó chưa đại diện cho mô hình cuối. Nhóm có thể dùng nó để \emph{hiệu
chuẩn rubric và tập chấm thử}, nhưng 500 cặp chính thức nên lấy từ checkpoint huấn
luyện trên tập chính.
\item \textbf{Chặn về khối lượng người:} yêu cầu ``hai người đánh giá độc lập'' có
nghĩa là \emph{hai con người}, mỗi người chấm toàn bộ số cặp. Phân tích tải nội bộ
của nhóm (dùng giả định 2 phút/cặp, không phải số đo) cho thấy 700 cặp thô cần
khoảng 23 giờ cho mỗi người chấm. Đây là một \emph{sàn thời gian}, không phải vấn
đề hiệu suất.
\end{enumerate}

\textbf{Điều nhóm không làm.} Không sinh nhãn bằng AI để thay người chấm. Nếu làm
vậy, con số ``mức đồng thuận'' sẽ không đo gì cả, và toàn bộ D.4 cùng D.5 (dựa lên
D.4) mất giá trị. Quy tắc nội bộ ghi rõ: không bịa điểm người chấm hay mức đồng
thuận.

\subsubsection{Hệ quả cho cầu nối A$\leftrightarrow$D (G5)}

Đề yêu cầu ít nhất một cài đặt Phần~A được dùng hoặc so sánh trực tiếp trong
Phần~D. Thiết kế của nhóm là đối chiếu \texttt{src/dpo.py} (A.4) với cách tính
chuẩn trên \emph{một minibatch dữ liệu ưu tiên thật} của D.4, so cả mất mát và
gradient. Phép đối chiếu này \textbf{chưa} thực hiện được vì chưa có dữ liệu đó.
Đây là hạng mục còn thiếu rõ ràng nhất của báo cáo.

%--------------------------------------------------------------------
\subsection{D.5 --- Đánh giá ba mốc mô hình}
\label{sec:d5}

\textbf{Yêu cầu.} So sánh base, SFT, SFT+DPO về độ chính xác đối chiếu nguồn,
hallucination, và từ chối khi ngoài phạm vi; blind human pairwise evaluation với
win rate kèm khoảng tin cậy; nêu giới hạn khi kiến thức thay đổi nhanh và hướng
RAG.

\textbf{Trạng thái: \chualam.}

\subsubsection{Giao thức dự kiến}

\begin{itemize}[leftmargin=2em]
\item Ba mốc model đánh giá trên \emph{cùng} tập test (tập khảo sát người thật của
D.2).
\item Ba nhóm chỉ số có định nghĩa trước: (i) đáp án nêu đúng điều kiện áp dụng và
dẫn đúng nguồn; (ii) tỷ lệ phát biểu không có căn cứ trong nguồn đã dẫn;
(iii) tỷ lệ từ chối đúng khi câu hỏi ngoài phạm vi hoặc chưa đủ căn cứ.
\item \textbf{Người đánh giá không biết đáp án thuộc model nào} (blind), thứ tự
trình bày được xáo.
\item Win rate báo cáo kèm khoảng tin cậy, ghi rõ đơn vị lấy mẫu.
\end{itemize}

\subsubsection{Vì sao chỉ số thứ ba quan trọng nhất}

Chỉ số (iii) đo đúng điều kiện thứ ba trong đặc tả sản phẩm
(Mục~\ref{sec:gioi-thieu}). Một mô hình đạt điểm cao ở (i) nhưng thấp ở (iii) là mô
hình nguy hiểm: nó trả lời tự tin cả khi không có căn cứ. Giả thuyết của nhóm là
DPO giúp nhiều nhất ở (iii), vì ``từ chối đúng cách'' là thứ khó dạy bằng SFT ---
dữ liệu SFT gồm các cặp câu hỏi--đáp án, nên nó dạy mô hình \emph{luôn} trả lời.

Đây là \emph{giả thuyết}, không phải kết quả. Nhóm không hứa DPO tốt hơn SFT và sẽ
báo cả kết quả âm nếu có.

\subsubsection{Phụ thuộc và giới hạn}

D.5 cần cả ba mốc model (phụ thuộc D.3 và D.4), tập test khảo sát thật (phụ thuộc
D.2), và người đánh giá mù --- là người thứ ba ngoài nhóm hoặc thành viên không
tham gia tạo nhãn. Không có đủ các điều kiện này thì bảng ba hệ không thể lập.

Về RAG: Đề tài 11 trong danh mục của đề là hướng truy xuất tăng cường. Nhóm chỉ
\emph{thảo luận} đây là hướng mở rộng hợp lý cho vấn đề kiến thức thay đổi nhanh
(xem hai mâu thuẫn ở Mục~\ref{sec:d1}), và \textbf{không} trình bày RAG như một
hạng mục đã triển khai.

%--------------------------------------------------------------------
\subsection{Tổng hợp trạng thái Phần D}

\begin{center}
\begin{tabular}{llp{0.46\textwidth}}
\toprule
\textbf{Mã} & \textbf{Trạng thái} & \textbf{Khoảng cách còn lại}\\
\midrule
D.1 & \motphan & Chưa phát hành bộ nguồn; còn kiểm kênh/phiên bản tại Hà Nội\\
D.2 & \motphan & 400/2.000 QA ứng viên; 24/400 kiểm thủ công; chưa có test khảo sát thật\\
D.3 & \motphan & Pilot 24 QA đã chạy (PPL validation 3,46; peak 2,402 GiB); chưa chọn cấu hình, chưa chạy tập chính, chưa đủ 5 seed\\
D.4 & \chualam & Đã có adapter pilot để sinh cặp; cần hai người chấm toàn bộ $\ge500$ cặp\\
D.5 & \chualam & Bị chặn bởi D.4 và tập test khảo sát; cần người đánh giá mù\\
\midrule
G5 & \chualam & Chưa đối chiếu A.4 trên minibatch ưu tiên thật của D\\
\bottomrule
\end{tabular}
\end{center}

\textbf{Đánh giá trung thực.} Phần~D chưa đạt, nhưng đường ống đã chạy hết một
vòng ở quy mô pilot: dữ liệu có dẫn nguồn $\to$ cổng phát hành $\to$ SFT trên GPU
$\to$ adapter sinh được đáp án. Điều này chứng minh thiết kế hoạt động.

Phần còn lại bị chặn bởi hai thứ khác nhau. Thứ nhất là \emph{khối lượng}: 394 QA
train cần người duyệt, và tập chính cần huấn luyện lại với nhiều cấu hình và $\ge5$
seed. Thứ hai, và không nén được bằng công sức của mã, là \emph{người thật}: người
khảo sát có văn bản đồng ý cho tập test (D.2), hai người chấm độc lập cho
$\ge500$ cặp (D.4), và người đánh giá mù (D.5). Nhóm chọn báo cáo đúng trạng thái
này thay vì điền số liệu không có nguồn.

\clearpage
\section{Tái lập}
\label{sec:tai-lap}

Đề yêu cầu: cố định seed; báo cáo trung bình $\pm$ độ lệch chuẩn trên ít nhất 5 lần
chạy độc lập hoặc 5 fold; siêu tham số chỉ chọn trên dữ liệu kiểm định
\cite[tr.~2, mục II]{de}. Mục này ghi điều kiện tái lập và nêu rõ phần nào còn
thiếu bằng chứng.

\subsection{Cấu trúc kho mã}

\begin{center}
\begin{tabular}{lp{0.60\textwidth}}
\toprule
\textbf{Đường dẫn} & \textbf{Nội dung}\\
\midrule
\texttt{src/} & Mã dùng lại: \texttt{attention.py} (A.1), \texttt{char\_lm.py}
(A.2), \texttt{lora.py} (A.3), \texttt{dpo.py} (A.4), \texttt{d3\_data.py}\\
\texttt{notebooks/} & Notebook điều khiển và giải thích, chạy tuần tự trong phiên
sạch\\
\texttt{tests/} & Bộ kiểm tự động\\
\texttt{scripts/} & Chuẩn bị dữ liệu D.2, đóng gói và huấn luyện D.3\\
\texttt{docs/theory/} & Phần B biên dịch độc lập\\
\texttt{docs/report/} & Báo cáo này\\
\texttt{docs/data/} & Nhật ký nguồn, quyền, quy trình riêng tư, datasheet\\
\texttt{docs/progress/} & Ma trận yêu cầu, bối cảnh, kế hoạch, nhật ký AI\\
\texttt{references/} & Năm bài báo và PDF đề gốc\\
\bottomrule
\end{tabular}
\end{center}

\subsection{Môi trường}

\textbf{Môi trường CPU cục bộ} (dùng cho toàn bộ phép kiểm A.1, A.3, A.4 và C.1--C.4):

\begin{center}
\begin{tabular}{ll}
\toprule
Python & 3.14.3\\
PyTorch & 2.12.1+cpu\\
CUDA khả dụng & \texttt{False}\\
\bottomrule
\end{tabular}
\end{center}

Bản PyTorch cài cục bộ là bản CPU, nên mọi con số ở các mục trên là số CPU. Máy có
GPU rời nhưng bản PyTorch hiện tại không dùng được nó; nhóm \emph{không} cài lại
PyTorch tùy tiện vì điều đó làm mất khả năng tái lập các kết quả đã ghi.

\textbf{Môi trường GPU cho D.3} khai báo ở \texttt{requirements-d3.txt}:
\texttt{transformers==4.57.6}, \texttt{peft==0.18.0},
\texttt{accelerate==1.12.0}, \texttt{bitsandbytes==0.48.2}. PyTorch và CUDA lấy từ
image Kaggle và \emph{không} thay thế, để tránh lệch giữa bản đã kiểm và bản chạy.

\textbf{Môi trường A.2} là một phiên Kaggle. Phiên hiển thị cấu hình
\texttt{GPU T4 x2} và trạng thái \emph{Version \#1 --- Successful}.

\subsection{Bộ kiểm tự động \dalam}

Chạy từ gốc kho mã:

\begin{quote}\small\ttfamily
python -m pytest -q
\end{quote}

Kết quả lần chạy ghi trong báo cáo này: \textbf{30 test qua, 1 test bỏ qua, 13
subtest qua} trong $8{,}38$ giây trên môi trường CPU nêu trên. Test bị bỏ qua là
test cần phụ thuộc không có trong môi trường cục bộ.

\subsection{Seed và số lần chạy}

\begin{center}
\begin{tabular}{lll}
\toprule
\textbf{Hạng mục} & \textbf{Seed} & \textbf{Số lần lặp}\\
\midrule
A.2 huấn luyện & 41, 42, 43, 44, 45 & 5 lần huấn luyện độc lập\\
C.1 đầu vào minh họa & \texttt{577893521} & 1 (bộ xác định, không lấy mẫu lại)\\
C.4 đầu vào minh họa & --- & 3 dấu margin $\times$ 4 giá trị $\beta$ (xác định)\\
\bottomrule
\end{tabular}
\end{center}

\textbf{Đơn vị lặp của A.2} là \emph{một lần huấn luyện hoàn chỉnh}, không phải một
lần sinh từ cùng checkpoint. Độ lệch chuẩn báo cáo là SD mẫu (mẫu số $n-1=4$) trên
5 giá trị PPL. Đây là độ biến thiên do khởi tạo và thứ tự trộn \emph{trên cùng một
split}; nó không phải khoảng tin cậy và không phải đánh giá trên tập test độc lập.

\subsection{Tách dữ liệu và quy tắc chọn siêu tham số}

\begin{itemize}[leftmargin=2em]
\item \textbf{A.2:} tách train (prefix 90\%) và validation (suffix 10\%) theo code
point \emph{trước} khi học từ vựng và \emph{trước} khi tạo cửa sổ. Từ vựng chỉ học
từ train. Phần đuôi không đủ một cửa sổ bị loại.
\item \textbf{D.2:} tách train/validation/test, kèm tách nhóm gần trùng (cùng họ
thủ tục, cùng câu hỏi). Tập test là \texttt{USER\_TEST}, có công cụ kiểm riêng để
chặn rò rỉ sang train.
\item \textbf{Quy tắc chung:} không dùng tập test để chọn cấu hình. Siêu tham số
của D.3 sẽ chọn trên validation; kế hoạch so $r$, $\alpha$, LR, epoch đã viết theo
quy tắc đó.
\end{itemize}

\subsection{Những gì còn thiếu để tái lập đầy đủ \chualam}

Nhóm liệt kê thẳng, thay vì tuyên bố đã đạt:

\begin{enumerate}[leftmargin=2em]
\item \textbf{Log môi trường và checkpoint của A.2.} Tệp \texttt{summary.json}
nhận được không chứa phiên bản PyTorch/CUDA và không kèm checkpoint. Nhóm chưa đối
chiếu được số GPU thực dùng. Việc cần làm: lưu log và output của phiên Kaggle đó.
\item \textbf{Hiện vật pilot D.2 không có trên máy hiện tại.} Các bản lưu
(\texttt{data/review/}, \texttt{.cache/}) nằm ngoài Git theo quy tắc riêng tư, nên
chỉ còn báo cáo JSON đã commit. Người tái lập cần nhận bàn giao riêng.
\item \textbf{Toàn bộ số đo của D.3--D.5} chưa tồn tại, nên chưa có gì để tái lập.
\item \textbf{Yêu cầu 5 lần chạy chưa áp dụng cho Phần D.} Kết quả 5 seed hiện chỉ
thuộc A.2. Không được suy ra rằng yêu cầu này đã đạt cho phần D.
\end{enumerate}

\subsection{Cách chạy lại từng phần}

\begin{center}
\begin{tabular}{ll}
\toprule
\textbf{Mục} & \textbf{Lệnh hoặc notebook}\\
\midrule
A.1 & \texttt{notebooks/01\_attention\_decoder.ipynb}\\
A.4 & \texttt{notebooks/02\_dpo\_loss.ipynb}\\
A.3 & \texttt{notebooks/03\_lora\_linear.ipynb}\\
C.2 & \texttt{notebooks/05\_transformer\_block\_cost.ipynb}\\
C.3 & \texttt{notebooks/06\_lora\_qlora\_cost.ipynb}\\
C.4 & \texttt{notebooks/07\_dpo\_loss\_gradient.ipynb}\\
C.1 & \texttt{notebooks/08\_causal\_attention\_example.ipynb}\\
D.3 & \texttt{notebooks/09\_d3\_sft\_kaggle.ipynb} (cần release D.2)\\
Phần B & \texttt{tectonic docs/theory/part-b.tex}\\
Báo cáo & \texttt{tectonic docs/report/report.tex}\\
\bottomrule
\end{tabular}
\end{center}

Các notebook C.1--C.4 không cần GPU, mạng, trọng số hay dữ liệu QA; chúng chạy tuần
tự trên CPU và tự tìm gốc kho mã.

\clearpage
\section{Đạo đức và pháp lý}
\label{sec:dao-duc}

Đề yêu cầu tuân thủ điều khoản nguồn, \texttt{robots.txt}, giới hạn truy cập và quy
định dữ liệu cá nhân; dữ liệu có yếu tố con người cần văn bản đồng ý và ẩn danh
trước khi lưu \cite[tr.~2, mục II]{de}. Mục này ghi những gì nhóm đã làm và những
ranh giới nhóm tự đặt.

\subsection{Nguồn và quyền sử dụng}

\subsubsection{Corpus của A.2}

Bản \emph{Truyện Kiều} Trương Vĩnh Ký 1911 thuộc phạm vi công cộng theo thời hạn
bảo hộ. Nhóm chuyển sang bản này \emph{sau khi} rà soát cho thấy căn cứ quyền đối
với một bản số hóa khác (revision \texttt{206607}) chưa đủ --- việc công khai để
đọc không đồng nghĩa được phân phối lại. Kết luận ``chưa đủ căn cứ'' được giữ lại
trong tài liệu thay vì xóa, để người đọc thấy đường ra quyết định.

Trong quá trình audit kỹ thuật, nhóm phát hiện \texttt{robots.txt} chặn đường
\texttt{/w/} của nguồn đó. Sau khi phát hiện, nhóm \emph{không} lặp lại truy cập
API hay raw và chỉ dùng route \texttt{/wiki/} được cho phép. Giới hạn này được ghi
rõ, không che giấu.

\subsubsection{Nguồn thủ tục của Phần D}

Nhóm phân loại nội dung theo quyền trước khi đưa vào dữ liệu
(Mục~\ref{sec:d1}): chính văn quy phạm và văn bản hành chính không thuộc phạm vi
bảo hộ quyền tác giả, nên có căn cứ dùng làm nguồn soạn QA kèm dẫn chứng đầy đủ;
còn trang tin, bài hướng dẫn và video thì chỉ dùng để \emph{tìm} và đối chiếu văn
bản, không chuyển nguyên trạng vào tập train.

Về truy cập: tuân thủ \texttt{Crawl-delay: 5} của cổng Hà Nội, truy xuất tuần tự,
không crawl toàn cổng, không truy cập các đường dẫn hồ sơ cá nhân, nộp/sửa hồ sơ,
SSO hay đăng nhập mà \texttt{robots.txt} chặn. Khi một host trả HTTP 404 cho
\texttt{robots.txt}, nhóm ghi rõ rằng \emph{không có tệp robots không phải là giấy
phép}.

Có một cơ chế trong pháp luật Việt Nam cho phép dùng có điều kiện dữ liệu được bảo
hộ cho nghiên cứu và huấn luyện AI (khoản 5 Điều 7 Luật SHTT; Điều 37a--37c
NĐ 17/2023 sửa bởi 134/2026). Nhóm ghi nhận cơ chế này nhưng \textbf{không} dựa vào
nó để duyệt các hiện vật pilot: nó không phải giấy phép cào hay phân phối snapshot.

\subsection{Dữ liệu có yếu tố con người}

Hai nhóm dữ liệu người xuất hiện trong Phần~D: câu hỏi khảo sát cho tập test của
D.2, và nhãn ưu tiên của D.4.

\subsubsection{Quy trình đồng ý và ẩn danh}

\begin{itemize}[leftmargin=2em]
\item Thu dữ liệu người cần \textbf{văn bản đồng ý} trước khi thu.
\item Dữ liệu được \textbf{ẩn danh trước khi lưu} làm dữ liệu nghiên cứu.
\item \textbf{Phiếu đồng ý có định danh để riêng, không đưa vào Git.} Kho mã có
quy tắc bỏ qua \texttt{/consent/} và \texttt{/surveys\_raw/}.
\item Có mốc để người tham gia \emph{rút} dữ liệu.
\end{itemize}

\subsubsection{Phân quyền trong nhóm}

Nhóm tách quyền truy cập theo vai trò, không theo mức tin cậy:

\begin{itemize}[leftmargin=2em]
\item Cao Văn Hà phụ trách lưu trữ vận hành: kho riêng mã hóa và sổ token liên kết
định danh.
\item Hà Mạnh Thế và Trần Duy Hưng \textbf{không} nhận phiếu đồng ý, bảng liên kết
định danh, hay khảo sát thô. Họ chỉ nhận dữ liệu DPO đã được phân quyền, và
\textbf{không} xem tập \texttt{USER\_TEST} trước khi giao thức đánh giá yêu cầu.
\item Hai người này phải \emph{xác nhận riêng bằng văn bản} rằng họ chỉ nhận dữ
liệu DPO được phân quyền.
\end{itemize}

Lý do: người chấm nhãn không cần biết ai đã trả lời khảo sát để làm việc của mình.
Giảm phạm vi truy cập giảm rủi ro, và việc xem trước tập test sẽ phá phép đo ở D.5.

\textbf{Trạng thái \chualam:} chưa có người tham gia nào được mời, chưa thu phiếu
đồng ý nào, và chưa có nhãn ưu tiên nào được tạo. Bộ công cụ đã soạn, hoạt động
thật thì chưa diễn ra.

\subsection{Ranh giới nhóm tự đặt về AI tạo sinh}

Ngoài các ràng buộc của đề, nhóm tự đặt bốn ranh giới và ghi chúng thành quy tắc
nội bộ:

\begin{enumerate}[leftmargin=2em]
\item \textbf{Không dùng AI tạo sinh thay khảo sát người thật.} Câu hỏi test phải
đến từ người dùng thật, vì mục đích của tập test là đo khả năng tổng quát sang cách
người dân diễn đạt (Mục~\ref{sec:d2}).
\item \textbf{Không bịa} nguồn, QA đã kiểm chứng, khảo sát người thật, điểm người
chấm, kết quả huấn luyện hay mức đồng thuận.
\item \textbf{Nhận xét của AI không được tính vào quota kiểm thủ công.} Trong vòng
rà pilot, AI có đối chiếu QA với đoạn trích và đề xuất sửa từng câu, nhưng số QA
``được người duyệt'' vẫn giữ bằng $0$ cho tới khi người thật điền phiếu.
\item \textbf{Không điền vào cột kết luận của người kiểm.} Phiếu kiểm có cột mẫu
minh họa, nhưng cột kết luận của người kiểm để trống.
\end{enumerate}

Ranh giới thứ ba và thứ tư là lý do báo cáo này ghi ``0 QA được người duyệt'' dù đã
có 20 QA qua kiểm cấu trúc tự động.

\subsection{Những gì không đưa vào kho mã}

Theo quy tắc riêng tư và quy định của đề, kho mã \emph{không} chứa: secret và khóa
truy cập; dữ liệu cá nhân, khảo sát thô, phiếu đồng ý; checkpoint và mô hình nền.
Hệ quả là một số hiện vật pilot chỉ còn báo cáo JSON đã commit, và người tái lập
cần nhận bàn giao riêng (Mục~\ref{sec:tai-lap}).

Về năm bài báo trong \texttt{references/}: chúng đã được theo dõi trong Git từ
trước. Quyền phân phối lại chưa được kiểm chứng, nên nhóm ghi nhận đây là rủi ro và
xin quyết định trước khi công khai kho mã, thay vì tự xóa hoặc tự công khai.

\subsection{Giới hạn về tính chính xác của nội dung pháp lý}

Sản phẩm là trợ lý tra cứu, không phải nguồn pháp lý. Hai điểm được ghi vào đặc tả
đầu ra ngay từ đầu (Mục~\ref{sec:gioi-thieu}):

\begin{enumerate}[leftmargin=2em]
\item Câu trả lời phải \textbf{dẫn nguồn}, vì giá trị pháp lý thuộc về văn bản gốc.
\item Khi ngoài phạm vi hoặc chưa đủ căn cứ, trợ lý phải \textbf{nói rõ} thay vì
đoán.
\end{enumerate}

Hai mâu thuẫn tìm được ở Mục~\ref{sec:d1} --- thời hạn cấp trích lục giữa hai quyết
định, và việc trộn biểu phí chứng thực với biểu lệ phí hộ tịch --- cho thấy rủi ro
này không phải lý thuyết. Một hệ thống trả lời tự tin về thời hạn ``3 ngày'' mà
không nêu văn bản căn cứ sẽ gây hại thật cho người dùng.

\clearpage
\section{Giới hạn và hướng mở rộng}
\label{sec:gioi-han}

\subsection{Giới hạn của các kết quả đã có}

\subsubsection{Phần A}

\begin{itemize}[leftmargin=2em]
\item \textbf{PPL của A.2 không so sánh được ra ngoài.} Giá trị $10{,}1948$ gắn với
một corpus cụ thể, một từ vựng ký tự 133 phần tử, một split cố định và một quy ước
chấm token. Nó không so được với PPL của mô hình khác dùng tokenizer khác. So sánh
perplexity chỉ có nghĩa khi cùng tokenizer, cùng dữ liệu và cùng quy ước chấm
(Mục~\ref{sec:b6-agg}).
\item \textbf{Độ lệch chuẩn $0{,}0628$ đo biến thiên hẹp.} Nó đo biến thiên giữa
các lần huấn luyện \emph{trên cùng một split}. Nó không bao gồm biến thiên do chọn
split khác, và không phải khoảng tin cậy.
\item \textbf{Chưa đối chiếu môi trường của A.2.} Thiếu log PyTorch/CUDA và
checkpoint, nên hạng mục tái lập của A.2 chưa đóng được
(Mục~\ref{sec:tai-lap}).
\item \textbf{Mô hình A.2 rất nhỏ} (2 khối, $\dmodel=64$). Nó là bằng chứng rằng
cài đặt đúng và huấn luyện được, không phải một mô hình ngôn ngữ hữu dụng.
\end{itemize}

\subsubsection{Phần B}

\begin{itemize}[leftmargin=2em]
\item \textbf{B.1 dựa trên một mô hình giả định.} Giả thiết độc lập, kỳ vọng $0$,
phương sai $1$ mô tả trạng thái tại khởi tạo, không phải phân phối của $q,k$ trong
mô hình đã huấn luyện. Kết quả giải thích lý do chọn thang điểm, không bảo đảm
không bão hòa.
\item \textbf{B.3 chỉ đúng với số chiều chẵn}, và Mục~\ref{sec:b3-odd} chứng minh
mệnh đề \emph{thất bại} ở chiều lẻ --- mà cài đặt của nhóm lại chấp nhận chiều lẻ.
Thí nghiệm hiện tại dùng chiều chẵn nên không chạm phải, nhưng đây là một lỗ hổng
thật giữa chứng minh và mã, không phải chi tiết hình thức.
\item \textbf{B.4 nói về chi phí số học, không nói về độ trễ.} Kết luận hợp nhất
không bảo đảm khớp từng bit trong số học dấu phẩy động, không tính chi phí tạo
$W'$, và không mở rộng thành kết luận về tốc độ khi đổi kernel hay tái lượng tử
hóa.
\item \textbf{B.5 chứng minh tính duy nhất của $\pi^*$, không phải của $\theta$.}
Việc mất mát DPO có nghiệm $\theta$ duy nhất hay không là câu hỏi khác, và không
được suy ra từ Mục~\ref{sec:b5}. Với miền đáp án không bị chặn reward, cần kiểm
riêng việc chuẩn hóa.
\end{itemize}

\subsubsection{Phần C}

\begin{itemize}[leftmargin=2em]
\item \textbf{C.1 và C.4 dùng đầu vào minh họa}, không phải bộ số giảng viên cấp.
Nhóm đã ghi rõ nguồn và cách tái lập, nhưng nếu nhận được đầu vào chính thức thì
phải thay số.
\item \textbf{C.2 chỉ đếm MAC của các tích ma trận} trên \emph{một} khối,
\emph{một} chuỗi, \emph{forward}. Không phải FLOP huấn luyện toàn mô hình. Các thao
tác không phải matmul (LayerNorm, softmax, ReLU, cộng bias) được nêu bằng
$\Theta(\cdot)$ chứ không cộng vào tổng.
\item \textbf{C.3 là ước lượng bộ nhớ, không phải phép đo.} Chưa cộng tái tính
activation, buffer giải lượng tử, workspace GEMM, CUDA context, tensor tạm của
optimizer và phân mảnh allocator. \textbf{Không} được kết luận từ bảng rằng một GPU
cụ thể chạy vừa SFT hay DPO.
\end{itemize}

\subsubsection{Phần D}

Giới hạn lớn nhất của báo cáo nằm ở đây: Phần~D chưa có kết quả thực nghiệm. Chi
tiết ở Mục~\ref{sec:phan-d}; tóm lại là D.2 còn 24/2.000 QA, D.3 chưa huấn luyện,
D.4 và D.5 chưa bắt đầu, và cầu nối A$\leftrightarrow$D (G5) chưa đóng.

\subsection{Rủi ro thiết kế đã nhận diện}

\subsubsection{Kiến thức thủ tục thay đổi nhanh}

Đây là rủi ro cấu trúc, không phải lỗi triển khai. Hai mâu thuẫn tìm được ở
Mục~\ref{sec:d1} đã minh họa: trong vòng chưa đầy một năm, thời hạn cấp trích lục
được quy định khác nhau ở hai quyết định, và nhóm không tìm được văn bản triển khai
tại Hà Nội để kết luận văn bản nào đang áp dụng.

Một mô hình được tinh chỉnh \emph{đóng băng} kiến thức tại thời điểm huấn luyện.
Khi văn bản thay đổi, mô hình trả lời sai mà không có cách nào tự biết. Ba cách
giảm nhẹ nhóm áp dụng trong phạm vi hiện tại:

\begin{enumerate}[leftmargin=2em]
\item Buộc câu trả lời dẫn số hiệu văn bản và phiên bản, để người dùng tự thấy đáp
án gắn với văn bản nào.
\item Soạn QA gắn với \emph{từng văn bản và phạm vi của nó}, thay vì phát biểu
chung ``hiện hành''.
\item Dạy mô hình từ chối khi chưa đủ căn cứ, thay vì nội suy.
\end{enumerate}

\subsubsection{Rủi ro về nhãn ưu tiên}

Nhãn ưu tiên là tín hiệu học duy nhất của DPO, nên chất lượng nhãn quyết định chất
lượng mô hình. Ba rủi ro:

\begin{itemize}[leftmargin=2em]
\item \textbf{Hai người chấm có thể lệch hệ thống} (ví dụ một người khắt khe hơn về
trích dẫn). Vì vậy phải \emph{đo} mức đồng thuận, không chỉ phân xử.
\item \textbf{Phân xử trước khi khóa nhãn sẽ làm hỏng phép đo đồng thuận.} Quy
trình đặt bước khóa trước bước phân xử chính vì lý do này.
\item \textbf{Cặp phải cùng prompt.} Theo \eqref{eq:b5-invert}, hạng
$\beta\log Z(x)$ chỉ triệt tiêu khi hai đáp án thuộc cùng một câu hỏi. Ghép đáp án
của hai câu hỏi khác nhau làm sai mô hình xác suất, không chỉ làm dữ liệu kém.
\end{itemize}

\subsubsection{Rủi ro đánh giá}

Nếu người đánh giá biết đáp án nào thuộc mô hình nào, win rate mất ý nghĩa. Nếu
người chấm nhãn D.4 cũng là người đánh giá D.5 và đã xem tập test, phép đo bị nhiễm.
Đây là lý do nhóm tách quyền truy cập \texttt{USER\_TEST}
(Mục~\ref{sec:dao-duc}).

\subsection{Hướng mở rộng}

\subsubsection{Truy xuất tăng cường (RAG)}

Hướng tự nhiên nhất cho vấn đề kiến thức thay đổi nhanh là truy xuất tăng cường:
thay vì ghi nội dung thủ tục vào trọng số, cho mô hình truy xuất văn bản tại thời
điểm trả lời và sinh đáp án dựa trên đoạn đã truy xuất. Cách này có hai lợi thế
đúng với bài toán: cập nhật văn bản không cần huấn luyện lại, và trích dẫn trở thành
hệ quả tự nhiên của kiến trúc chứ không phải hành vi phải dạy.

Đây là Đề tài 11 trong danh mục của đề. Nhóm \textbf{chỉ thảo luận} đây là hướng mở
rộng và \textbf{không} trình bày RAG như hạng mục đã triển khai; mở rộng sang RAG
không nằm trong phạm vi Đề tài 9.

\subsubsection{Các mở rộng khác trong phạm vi Đề tài 9}

\begin{itemize}[leftmargin=2em]
\item \textbf{Vá B.3 trong mã:} buộc $\dmodel$ chẵn, hoặc bỏ tọa độ $\sin$ lẻ cặp,
để tính chất dịch vị trí đúng với mọi cấu hình mà mã cho phép.
\item \textbf{Mở rộng adapter ngoài Q/V:} C.3 tính cho Q/V theo đề. Gắn thêm
adapter vào K, O hoặc các ma trận FFN làm tăng số tham số học theo
\eqref{eq:b4-count} và có thể đổi kết quả; đây là câu hỏi thực nghiệm chưa trả lời.
\item \textbf{Đo peak VRAM thật} để thay ước lượng C.3 bằng số đo, ghi cả peak
allocated và peak reserved.
\item \textbf{Mở rộng tập test khảo sát} sang nhiều nhóm người dùng, vì cách diễn
đạt khác nhau giữa các nhóm tuổi và mức quen thuộc với thủ tục hành chính.
\end{itemize}

\subsection{Nhóm không hứa điều gì}

Để tránh ngộ nhận, nhóm nêu rõ ba điều \emph{không} được suy ra từ báo cáo này:

\begin{enumerate}[leftmargin=2em]
\item \textbf{Không hứa DPO tốt hơn SFT.} Giả thuyết ở Mục~\ref{sec:d5} là giả
thuyết. Nhóm sẽ báo cả kết quả âm nếu đo được.
\item \textbf{Không coi bằng chứng kỹ thuật là bằng chứng quyền.} Việc audit chạy
được không có nghĩa nhóm có quyền huấn luyện hay phân phối dữ liệu đó.
\item \textbf{Không coi kế hoạch là kết quả.} Mọi hạng mục đã lập kế hoạch nhưng
chưa chạy đều mang nhãn \chualam.
\end{enumerate}

\clearpage
\appendix
\section{Phụ lục --- Khai báo phạm vi dùng AI tạo sinh}
\label{sec:phu-luc-ai}

Đề yêu cầu khai báo phạm vi dùng AI tạo sinh trong phụ lục; từng thành viên có thể
bị vấn đáp; không dùng bản AI miễn phí \cite[tr.~2, mục II]{de}. Phụ lục này khai
báo trung thực, kể cả những phần AI làm phần lớn công việc.

Nguồn chi tiết theo từng phiên là \texttt{docs/progress/ai-use.md}, được cập nhật
liên tục trong suốt dự án.

\subsection{Công cụ và trạng thái xác minh}

Công cụ: một trợ lý lập trình AI chạy trong workspace của nhóm, có hỗ trợ soạn tài
liệu.

\textbf{Điểm cần nêu rõ:} nhóm \textbf{chưa xác minh} loại tài khoản hay gói dịch
vụ AI đang dùng có phù hợp quy định ``không dùng bản AI miễn phí'' hay không. Nhóm
không diễn giải quy định này thành cấm dùng Colab hay Kaggle bản miễn phí, vì đó là
nền tảng tính toán, không phải công cụ AI tạo sinh. Việc xác minh tài khoản là hạng
mục còn mở.

\subsection{Phạm vi AI đã làm}

\subsubsection{Tài liệu, kế hoạch và hồ sơ}

AI kiểm kê Git và tài liệu; đọc và đối chiếu các PDF trong \texttt{references/};
soạn ma trận yêu cầu, đề cương, kế hoạch dữ liệu, hướng dẫn kho mã, bối cảnh bàn
giao, kế hoạch phối hợp nhóm và lịch chuẩn bị vấn đáp; tạo mẫu LaTeX, bộ lọc Pandoc
và script sinh/biên dịch PDF; tổ chức lại README.

\subsubsection{Mã Phần A và bộ kiểm}

\begin{itemize}[leftmargin=2em]
\item \textbf{A.1:} AI viết \texttt{src/attention.py}, notebook dẫn xuất và bộ
kiểm; thư viện chuẩn của PyTorch chỉ dùng ở phía đối chiếu. Chạy trên CPU cục bộ.
\item \textbf{A.2:} AI viết LM dùng decoder A.1, mã đọc dữ liệu có kiểm quyền, vòng
huấn luyện, tính PPL, lưu/khôi phục checkpoint, notebook và bộ kiểm. AI không chạy
Kaggle thay nhóm.
\item \textbf{A.3:} AI triển khai lớp LoRA, notebook và bộ kiểm CPU; chạy đối chiếu
số tham số và phép hợp nhất sau cập nhật khác $0$.
\item \textbf{A.4:} AI triển khai hàm DPO, notebook và bộ kiểm CPU bằng số minh
họa; đối chiếu mất mát và gradient.
\end{itemize}

AI cũng đã phát hiện, tái hiện và sửa một lỗi trong script corpus (theo redirect
trước khi kiểm URL đích), kèm kiểm thử hồi quy.

\subsubsection{Phần B và Phần C}

\begin{itemize}[leftmargin=2em]
\item \textbf{Phần B:} AI đối chiếu sáu yêu cầu với đề tr.~18, đọc các phần liên
quan của \cite{attention}, \cite{lora}, \cite{dpo}, và \textbf{soạn toàn bộ bản
nháp LaTeX B.1--B.6}: phát biểu, giả thiết, các bước chứng minh, diễn giải và ví dụ
minh họa. Ở một lượt rà sau đó, AI tìm ra lỗ hổng số chiều lẻ của B.3 và soạn phần
bổ sung ở Mục~\ref{sec:b3-odd}.
\item \textbf{C.1:} theo yêu cầu của nhóm, AI sinh ma trận Q/K/V với seed
\texttt{577893521}, lưu đầu vào tái lập, soạn phép tính tay từng token và đối chiếu
với mã A.1. Việc giảng viên không cung cấp ma trận là thông tin do nhóm báo.
\item \textbf{C.2:} AI dẫn xuất số tham số và MAC theo kích thước tensor của
\texttt{DecoderBlock}, rồi chạy notebook CPU đối chiếu \texttt{numel()} và FLOP do
profiler ước lượng.
\item \textbf{C.3:} AI đối chiếu cấu hình Qwen khóa revision, mã Qwen2 trong
Transformers và giả định double quantization của \cite{qlora}; viết notebook và
chạy trên CPU để kiểm số tham số bằng lớp A.3.
\item \textbf{C.4:} AI tự đặt bốn log-xác suất minh họa, trình bày phép thế số, và
kiểm mất mát/gradient với A.4 và autograd.
\end{itemize}

\subsubsection{Phần D}

\begin{itemize}[leftmargin=2em]
\item \textbf{D.1:} AI truy tìm văn bản công khai, đối chiếu số hiệu, ngày, phạm
vi, phụ lục và căn cứ quyền tác giả; lưu URL, thời điểm và SHA-256; kiểm chữ ký CMS
về toàn vẹn nội dung trên bốn PDF Hà Nội. AI không đăng nhập, không lấy hồ sơ công
dân, không tự phê duyệt nguồn.
\item \textbf{D.2:} AI soạn recipe từ chính văn, kiểm SHA-256 và đoạn trích, triển
khai CLI chọn mẫu phân tầng và cổng xuất SFT, cài công cụ kiểm tách
\texttt{USER\_TEST}. AI biên soạn 20 QA nghiên cứu và đối chiếu 21 đoạn trích.
\item \textbf{D.3:} AI chuẩn bị bộ nạp release D.2, mã QLoRA, notebook Kaggle,
script đóng gói và kế hoạch so cấu hình trên validation; chạy kiểm phần mềm trên
CPU.
\item AI soạn bộ tài liệu đồng ý, khảo sát, quy trình riêng tư và rubric/phiếu
chấm.
\end{itemize}

\subsubsection{Báo cáo này}

AI soạn cấu trúc và toàn bộ bản nháp của báo cáo này từ các hiện vật đã có trong
kho mã, kèm datasheet và model card. AI đối chiếu lại từng con số với notebook, tệp
tài liệu nguồn và ma trận yêu cầu; chạy bộ kiểm tự động và ghi kết quả thực tế. AI
không tạo kết quả thực nghiệm mới và không đổi trạng thái của hạng mục nào.

\subsection{Phạm vi AI \emph{không} làm}

Danh sách này quan trọng ngang danh sách trên:

\begin{enumerate}[leftmargin=2em]
\item \textbf{Không thực hiện khảo sát con người} và không tạo bản ghi khảo sát.
Chưa mời hay thu dữ liệu của bất kỳ người tham gia nào.
\item \textbf{Không gán nhãn ưu tiên thật} và không tạo mức đồng thuận.
\item \textbf{Không điền phiếu kiểm thủ công.} Trong vòng rà pilot, AI chạy
validator ngoại tuyến và lưu báo cáo riêng, nhưng số QA được người duyệt vẫn giữ
bằng $0$. Nhận xét của AI \emph{không} được tính vào quota kiểm thủ công $\ge20\%$.
\item \textbf{Không chạy huấn luyện thay nhóm} trên Colab hay Kaggle, không tạo
dataset/notebook trên tài khoản nền tảng của nhóm.
\item \textbf{Không tự phê duyệt quyền dùng nguồn} và không ký xác nhận thay người
có thẩm quyền. Kết quả audit kỹ thuật không được dùng làm attestation về quyền.
\item \textbf{Không thay bộ Q/K/V chính thức của Phần C} --- bộ minh họa luôn được
ghi nhãn là minh họa.
\item \textbf{Không tạo kết quả đánh giá} hay số đo GPU.
\end{enumerate}

\subsection{Dữ liệu đầu vào cho AI}

Yêu cầu của nhóm; sáu PDF có sẵn trong workspace (PDF đề gốc và năm bài báo); các
trang web công khai đã ghi URL và thời điểm trong nhật ký nguồn; và những thông số
phiên Kaggle/Colab \emph{do thành viên nhóm cung cấp}. Các thông số do thành viên
cung cấp được ghi là thông tin người dùng cung cấp, \emph{không} nhận là kết quả do
AI đo.

\subsection{Trách nhiệm của nhóm}

Những việc sau thuộc trách nhiệm thành viên và \emph{không} được chuyển cho AI:

\begin{itemize}[leftmargin=2em]
\item Đọc lại nguồn và xác minh các giả thiết.
\item Lựa chọn phạm vi dự án và mô hình.
\item \textbf{Hiểu và giải thích được toàn bộ mã và chứng minh khi vấn đáp.} Mỗi
thành viên tự giải lại các chứng minh của Phần~B và chạy lại các phép kiểm, không
dựa vào việc AI đã viết chúng.
\item Thực hiện kiểm định thủ công, khảo sát, gán nhãn ưu tiên và đánh giá mù.
\item Ghi nhận sửa đổi do thành viên thực hiện, công cụ/phiên bản và phạm vi AI ở
các phiên sau.
\end{itemize}

\textbf{Trạng thái tại thời điểm viết:} nhóm chưa ghi tên người kiểm duyệt nội dung
cuối vì chưa có xác nhận. Phần~B vẫn mang nhãn bản nháp chưa được nhóm duyệt
(Mục~\ref{sec:phan-b}). Văn bản do AI soạn \textbf{không} được coi là nguồn học
thuật thay cho năm bài báo và đề gốc.


%====================================================================
\clearpage
\begin{thebibliography}{9}
\bibitem{de}
Phạm Văn Khánh. \emph{Danh mục 20 bài tập lớn: Học máy cơ bản, mô hình ngôn ngữ lớn,
học tăng cường và tích hợp}. Bản ngày 12/11/2026, 46 trang; Đề tài 9, tr.~18--19.

\bibitem{attention}
Vaswani, A. và cộng sự. \emph{Attention Is All You Need}. NeurIPS 2017;
bản \href{https://arxiv.org/abs/1706.03762v7}{arXiv:1706.03762v7}.

\bibitem{lora}
Hu, E. J. và cộng sự. \emph{LoRA: Low-Rank Adaptation of Large Language Models}.
ICLR 2022; bản \href{https://arxiv.org/abs/2106.09685v2}{arXiv:2106.09685v2}.

\bibitem{qlora}
Dettmers, T. và cộng sự. \emph{QLoRA: Efficient Finetuning of Quantized LLMs}.
NeurIPS 2023; bản \href{https://arxiv.org/abs/2305.14314v1}{arXiv:2305.14314v1}.

\bibitem{rlhf}
Ouyang, L. và cộng sự. \emph{Training language models to follow instructions with
human feedback}. NeurIPS 2022; bản \href{https://arxiv.org/abs/2203.02155v1}{arXiv:2203.02155v1}.

\bibitem{dpo}
Rafailov, R. và cộng sự. \emph{Direct Preference Optimization: Your Language Model
Is Secretly a Reward Model}. NeurIPS 2023;
bản \href{https://arxiv.org/abs/2305.18290v3}{arXiv:2305.18290v3}.

\bibitem{qwen}
Qwen Team. \emph{Qwen2.5-1.5B-Instruct}. Model card trên Hugging Face, giấy phép
Apache-2.0; revision được khóa trong \texttt{scripts/} của kho mã.
\end{thebibliography}

\end{document}
